% Modernised reading — fonds Alexandre Grothendieck, shelfmark 2, pages 1-44.
%
% This is an INTERPRETATION, not a transcription. It re-reads the pages in
% today's notation and vocabulary, states the hypotheses the manuscript leaves
% implicit, and drops the critical apparatus so that the mathematics reads
% straight through. Everything the brackets and underlines carried moves into a
% footnote. It is held to being mathematically correct as it stands; where the
% manuscript is loose or mistaken, the true statement is what appears here and
% the footnote says what the page has. In any dispute of substance the
% transcription governs: cite batch-01/02/03.fr.tex.
%
% Pass: Opus 5 (claude-opus-5), 2026-09-19 — first pass, unchecked against the
% pages by a human.
%
% Scope: one document for the shelfmark, its three batches read as one folder.
% Twenty-eight of the forty-four leaves carry mathematics; five are covers in
% his hand, whose words become section headings here; the rest are blank or are
% a typescript that is not his.
%
% SPINE. The folder is not one text, and its unity is not chronological: it is
% a single NUMBER, (-1)^{i-1}(i-1)!, appearing in four disguises on leaves that
% cannot all have been written at one date.
%
%   as the top Chern class of a point in P^r      c(xi^r) = 1 + (-1)^{r-1}(r-1)! eta^r   (p. 9)
%   as the defect of the comparison CK -> GK      phi psi = (-1)^{i-1}(i-1)!             (p. 21)
%   as the reason a nonzero class can have        (r-1)! x' = 0  =>  c = 1               (p. 13)
%     trivial total Chern class
%   as the coefficient of log c in terms of ch    log c = sum (-1)^{m-1}(m-1)! ch_m      (p. 43)
%
% The fourth is not written on any leaf; it is what page 43's exp(...) is the
% first two terms of, and it is what makes the other three one fact. Stated in
% the body as this edition's own observation, never as the page's.
%
%   pp. 2-3    RIEMANN-ROCH ON A CURVE, BEFORE THE LETTER K. The ring E of
%              classes of coherent sheaves, its subgroups E_0, N, D, L and the
%              subalgebra Phi; E = Phi + D, E/N = Z + P; R.R. in one line,
%              int(F) = K_1(c(F)(1 - K/2)).
%   pp. 4-7    THE SAME IN K-THEORY, dim <= 1. K^.(X) and K_.(X), rank and
%              length augmentations, the gamma-filtration, Pic(X), pi_0(X), the
%              comparison c = ch into B^.(X), B_.(X), and the map K^. -> K_.
%              which is multiplication by the class of O_X.
%   pp. 9-13   "CONTRE-EXEMPLES". The total Chern class of i_! along
%              X -> X x P^r; the factor (r-1)!; the polynomials Q^{(r)}_i; and
%              a nonzero class whose total Chern class is 1.
%   pp. 15-18  MODULES ON A RINGED TOPOS CARRYING A GROUP ACTION. Restriction
%              and induction, the projection formula, the external product, the
%              induction product over the symmetric groups, and the formulas
%              for T^n. This is where the lambda-operations of pp. 20-29 come
%              from; no leaf says so.
%   pp. 20-29  "CAS DES SCHEMAS": the lambda-ring formalism. Special
%              lambda-rings with geometric elements, the universal ring B, the
%              Chern ring CK, the gamma-filtration and GK, phi_K : CK -> GK,
%              the factor (-1)^{i-1}(i-1)!, and the counterexample
%              K = Z[L]/(L-1)^n = K(P^{n-1}).
%   pp. 32-37  "CONSTRUCTION DES CHERN COVARIANTS". D_parf(f), K with supports
%              in a closed Z or a family Phi, the Chow groups A_Phi(X), the
%              Chern classes between them, and R.R. as the measure of the
%              defect of covariance.
%   pp. 40-44  "R.R. DES SURFACES". The GRR square, (1), (1 bis), (2), (1 ter),
%              the fibration, a curve in a curve, a covering, a curve in a
%              non-singular surface, and — one identification short of being
%              written — the adjunction formula.
%
% What only a folder-wide reading can say, and which is written in the body:
%   - the 1 - K/2 of page 3 IS the tau_h(X) of page 42: the Todd class of a
%     curve, five wrappers and (at least) one campaign apart, under two names.
%   - the K_1 of page 3 IS the K_p of page 41 and the K_2 of page 44: one
%     operator, the degree of the component of degree p.
%   - the (r-1)! of page 9 IS the (i-1)! of page 21. The geometric
%     counterexample of pages 9-13 and the algebraic one of page 21 are the
%     same phenomenon: torsion that CK sees and GK does not.
%   - pages 15-18 construct the operations that pages 20-29 axiomatise.
%   - the gamma-filtration is on page 4 as well as on page 20, so it does not
%     separate the folder's campaigns.
%
% NOTATION fixed for the folder. His own varies from run to run; one name is
% chosen here and the variants are footnoted where they occur.
%   ch          — the Chern character. He writes c = ch (pp. 4, 6), ch
%                 (pp. 36-37) and varphi (pp. 40-44).
%   Todd        — the Todd class. He writes 1 - K/2 (p. 3) and tau_h (pp. 40-44).
%   c^i, c      — Chern classes and the total Chern class; c-tilde when the rank
%                 is carried alongside, as the pair [eps(x), 1 + c^1 + ...].
%   A-tilde     — for a graded ring A, the set Z x (1 + A^{>=1}) with Whitney
%                 sum as its addition and the universal Q_n as its product.
%   deg         — the degree of the component of degree p. He writes K_1 (p. 3),
%                 K_p (p. 41), K_2 (p. 44).
%   k_X, k_Y    — canonical classes. He writes K (p. 3), k and k' (p. 42),
%                 k_x, k_y, k_Y, k_X (pp. 42-44).
%   V, V', V''  — vector bundles on pages 2-3, where E is the ring of classes
%                 and cannot also be a bundle. The leaf uses E for both.
%   zeta        — the generator of CK on page 21, where he writes xi, which is
%                 already 1 - L^{-1} on pages 9-13 and a K-class on page 42.
%   h           — the first Chern class of the twisting sheaf on page 42, where
%                 he writes xi, which on the same leaf is also the K-class (n,d).
%   E           — the ring of classes (pp. 2-3), a distinguished element E_s
%                 (pp. 27-29), a vector bundle (pp. 41-43). Three objects, one
%                 letter, in three runs; kept, and separated in the spine section.
%   X, Y        — on pages 42 to 44, X is the source curve and Y the target
%                 throughout, which is what the formulas require and what the
%                 prose of those leaves contradicts four times. Each is footnoted.
%
% SUBSTANTIVE DEPARTURES, each footnoted at the point of departure:
%   - p. 9 / p. 11: the exponent of (-1) is written r+1 on one leaf and r-1 on
%     the other. The two are equal. The transcription flags a variance that is
%     not one, and the reading says so.
%   - p. 9: the expansion of c-tilde(x) * (1 - i eta) is transcribed without the
%     binomial coefficients the formula carries. Restored here.
%   - p. 11: the leaf computes the leading coefficient of Q^{(r)}_i only for
%     i = 1. The general form its own method gives is stated here and marked as
%     supplied.
%   - p. 13: the last displayed member is read 0 where the total Chern class is
%     1. Both are said; the two differ only by what c-tilde denotes.
%   - p. 20: phi psi is written (-1)^i (i-1)!; page 21 proves (-1)^{i-1}(i-1)!,
%     which is what holds.
%   - p. 22: he strikes the line c^i(xy) = Q_i(...) and rewrites it with
%     c-tilde on the left. The struck line is the correct one.
%   - p. 18: w_{m,n} is written as a map out of S_m. Its source is the wreath
%     product S_n wr S_m; supplied here.
%   - p. 40: (1 bis) is correct with T_{Y/X} read as the normal bundle of X in
%     Y. Said once.
%   - p. 42: the left-hand side is written varphi_X where the class lives on Y.
%   - p. 43: f_*(O_X) is called a bundle on X where it is one on Y.
%   - p. 44: the surface is Y in the formulas and X in the sentence.
%
% WHAT THE FOLDER ANNOUNCES AND NEVER ESTABLISHES: the universal polynomial for
% gamma^i(xy) (p. 22 stops on it); Propositions 1 and 2 of page 20, of which
% only the statement of the third survives, through its proof on page 21; the
% proposition of page 21 marked "à prouver directement?" and struck through,
% taken up in another form on page 27; the three questions of page 34 (A_Phi(X)
% for general X, the lambda-procedure, cohomological Chern classes); the
% filtration on K^._X(X'), a title on page 37 with nothing under it; the shape
% of (2) on page 40; and the adjunction formula, which page 44 is one
% identification short of writing when the folder stops.

\documentclass[11pt,a4paper]{article}
% Le préambule commun à toutes les transcriptions du fonds Grothendieck.
%
% Il tient en une idée : **l'incertitude se compose, elle ne se résout pas.**
% Une transcription qui lisse un mot douteux a détruit la seule chose pour
% laquelle on la faisait — un lecteur doit pouvoir savoir, à chaque endroit,
% ce qui était sur le feuillet et ce qui a été ajouté. Les six macros qui
% suivent sont donc l'appareil critique, et non de la décoration.
%
% Les mêmes commandes sont reconnues par `scripts/render.mjs`, qui produit la
% vue de lecture du site. Une macro ajoutée ici doit l'être aussi là-bas,
% faute de quoi le rendu échouera bruyamment — ce qui est voulu : un rendu
% silencieusement faux est pire qu'un rendu absent.

\ProvidesPackage{grothendieck}[2026/08/08 Transcriptions du fonds Grothendieck]

\usepackage{amsmath,amssymb,amsthm}
\usepackage{mathrsfs}
\usepackage{stmaryrd}
% Les diagrammes commutatifs sont partout dans ce fonds : c'est la langue
% même de la géométrie algébrique de Grothendieck, et une transcription qui
% les rendrait en prose ne transcrirait rien.
\usepackage{tikz-cd}
% Français par défaut — c'est la langue des feuillets — mais l'anglais est
% chargé aussi : la traduction et le résumé sont des documents anglais, et la
% typographie française (espaces avant les deux-points) y serait une faute.
% Ces documents basculent par \selectlanguage{english} en tête de corps.
\usepackage[english,french]{babel}
\usepackage{fontspec}
\usepackage{geometry}
\geometry{a4paper,margin=25mm,marginparwidth=28mm}
\usepackage{marginnote}
\usepackage{xcolor}
% ulem, not soul. `soul`'s \st and \ul are built on a character-by-character
% scan that XeLaTeX's font machinery breaks: the document fails with an
% undefined control sequence rather than degrading. ulem does the same two jobs
% and is Unicode-engine safe. `normalem` stops it hijacking \emph, which the
% transcriptions use for Grothendieck's own emphasis.
\usepackage[normalem]{ulem}
\usepackage{hyperref}
\hypersetup{colorlinks=true,linkcolor=black,urlcolor=black,pdfborder={0 0 0}}

% --- Métadonnées du lot -----------------------------------------------------
% Reprises telles quelles par le script de rendu, qui les affiche en tête de
% la vue de lecture. Elles fixent ce que le fichier prétend couvrir : sans
% elles, une transcription flotte sans point d'attache dans un fonds de seize
% mille feuillets.

\newcommand{\folder}[1]{\gdef\grothFolder{#1}}
\newcommand{\batch}[1]{\gdef\grothBatch{#1}}
\newcommand{\leaves}[2]{\gdef\grothFirst{#1}\gdef\grothLast{#2}}
\newcommand{\foldertitle}[1]{\gdef\grothFoldertitle{#1}}
\gdef\grothFolder{?}\gdef\grothBatch{?}\gdef\grothFirst{?}\gdef\grothLast{?}\gdef\grothFoldertitle{}

% --- Le résumé --------------------------------------------------------------
%
% Il ouvre la lecture modernisée, sous le titre « Résumé », et il est composé
% en petit. Ce n'est pas un ornement : c'est le seul passage du document qui
% ne soit pas des mathématiques, et un lecteur doit pouvoir le reconnaître
% avant de l'avoir lu — puis le sauter s'il connaît déjà le sujet, ou s'y
% arrêter s'il ne le connaît pas. Le filet qui le suit marque la reprise du
% corps.

\newenvironment{resume}
  {\small\setlength{\parskip}{0.45em}}
  {\par\medskip\hrule\medskip}

% --- Mots-clés ---------------------------------------------------------------
%
% En anglais, en clôture du résumé de la lecture modernisée : le vocabulaire
% moderne sous lequel chercher ce que les feuillets construisent. Le manifeste
% du site (`npm run manifest`) les extrait pour étiqueter la cote entière —
% cette ligne est la source unique des tags, il n'y a pas de fichier à côté.
\newcommand{\keywords}[1]{\par\smallskip\noindent
  {\footnotesize\textsc{Keywords} --- \textit{#1}}\par}

% --- L'appareil critique ----------------------------------------------------

% Le numéro de feuillet, en marge. C'est l'ancre qui permet de régler un
% désaccord sur une formule en regardant *un* feuillet plutôt qu'une chemise
% de six cents. Le site s'en sert aussi pour tourner les pages du fac-similé
% quand on fait défiler la transcription.
\newcommand{\leaf}[1]{%
  \marginnote{\footnotesize\color{gray!70}\textsf{#1}}%
  \label{leaf:#1}\ignorespaces}

% Illisible : on ne devine pas. Un mot inventé qui se lit comme les autres est
% le pire résultat possible.
\newcommand{\ill}{\textcolor{red!60!black}{[\,\ldots\,]}}

% Lecture douteuse : le mot est donné, et signalé comme incertain.
\newcommand{\uncertain}[1]{\uline{#1}}

% Ajout de l'éditeur — une lettre manquante, un mot sous-entendu.
\newcommand{\add}[1]{\textcolor{blue!50!black}{[#1]}}

% Biffé par Grothendieck lui-même. Ce qu'il a rayé fait partie du document :
% une rature dit souvent où la pensée a changé de direction.
\newcommand{\struck}[1]{\sout{#1}}

% Note du transcripteur, sur l'état matériel du feuillet.
\newcommand{\note}[1]{\par\smallskip{\small\color{gray!60!black}\textit{[#1]}}\par\smallskip}

% Note marginale de Grothendieck, distincte du corps du texte.
\newcommand{\marginal}[1]{\par\smallskip\noindent\hspace{1em}%
  {\small\color{gray!60!black}#1}\par\smallskip}

% --- Titre ------------------------------------------------------------------

\AtBeginDocument{%
  \begin{center}
    {\large\bfseries Fonds Alexandre Grothendieck --- cote n\textsuperscript{o}~\grothFolder}\\[2pt]
    {\itshape\grothFoldertitle}\\[4pt]
    {\small feuillets \grothFirst--\grothLast\ (lot \grothBatch)}\\[2pt]
    {\footnotesize\color{gray!60!black}Transcription établie d'après le fac-similé numérisé
     par l'Université de Montpellier.\\
     Les lectures incertaines sont soulignées, les illisibles notées [\,\ldots\,],
     les ajouts entre crochets.}
  \end{center}
  \vspace{1em}}

\endinput


\folder{2}
\pages{1}{44}
\dating{[1957]}
\watermark{Édition de démonstration\\interprétation personnelle de l'œuvre}
\foldertitle{Théorème de Riemann-Roch / RR [Riemann-Roch] Notes Antiques (antérieur[es à] 1958) : notes manuscrites (s.d.). — lecture modernisée du dossier entier}

\begin{document}

\section*{Résumé}

\begin{resume}
Ces quarante-quatre feuillets tournent autour d'un théorème et autour d'un
nombre. Le théorème est celui de Riemann-Roch. Le nombre est
$(-1)^{i-1}(i-1)!$, et c'est lui, plutôt que le théorème, qui donne au dossier
son unité.

Voici d'abord le théorème, dans sa forme la plus ancienne. Prenons une courbe
algébrique compacte et lisse --- on peut se la représenter comme une surface :
la sphère, le tore, le bretzel à $g$ trous. Donnons-nous dessus un nombre fini
de points et, en chacun, un ordre de pôle autorisé. Combien existe-t-il de
fonctions sur la courbe n'ayant de pôles qu'aux endroits permis, et pas plus
forts que permis ? Ces fonctions forment un espace vectoriel de dimension
finie, et la question est de calculer cette dimension. Il se trouve qu'elle ne
dépend presque que de deux entiers : le nombre total de pôles autorisés, que
l'on appelle le \emph{degré}, et le nombre $g$ de trous, que l'on appelle le
\emph{genre}. Presque : la formule exacte fait apparaître un terme correctif,
qui est lui-même une dimension du même type pour une autre donnée de pôles.
C'est ce que l'on appelle le théorème de Riemann-Roch, et sous cette forme il
date du XIX\textsuperscript{e} siècle.

Le terme correctif est gênant tant qu'on le prend pour un reste. Il devient
naturel dès qu'on cesse de compter une dimension pour compter une
\emph{différence alternée} de dimensions --- première dimension moins seconde.
Cette différence porte un nom, la caractéristique d'Euler-Poincaré, et elle a
sur chacun des deux termes pris isolément un avantage décisif : elle est
\emph{additive}. Si un objet se décompose en deux morceaux, sa caractéristique
est la somme des leurs. Or c'est exactement ce dont on a besoin pour calculer :
une quantité additive se laisse suivre à travers les découpages.

De là vient l'idée qui organise tout le dossier. Considérons tous les objets
que l'on peut poser sur la courbe --- ce sont les \emph{faisceaux cohérents},
mais l'essentiel est qu'il y en a beaucoup et qu'ils se découpent les uns dans
les autres --- et décidons de ne retenir de chacun qu'une étiquette, sa
\emph{classe}, avec une seule règle : quand un objet se découpe en deux, sa
classe est la somme des deux autres. On obtient ainsi un groupe, et même, le
produit tensoriel fournissant une multiplication, un anneau. Cet anneau ne se
souvient de rien d'autre que de ce qui est additif ; c'est toute sa force. On
l'appelle aujourd'hui le \emph{groupe de Grothendieck}, ou l'anneau
$K$ ; ces pages l'écrivent tantôt $E$, tantôt $K$, ce qui est déjà un indice
sur leur date. Le théorème de Riemann-Roch, une fois ces classes en main, tient
littéralement sur une ligne, et la page 3 l'écrit sur une ligne.

Reste à savoir mesurer une classe. Il y a pour cela deux règles graduées, et la
tension entre elles est le sujet véritable du dossier. La première, les
\emph{classes de Chern}, est à coefficients entiers : elle ne ment pas, mais
elle se comporte mal vis-à-vis de l'addition et de la multiplication, où
apparaissent des formules compliquées. La seconde, le \emph{caractère de
Chern}, est parfaitement additive et multiplicative --- mais il faut, pour
l'écrire, autoriser des dénominateurs, c'est-à-dire sortir des entiers. On veut
évidemment les deux à la fois : la bonne algèbre et les coefficients entiers.
On ne les a pas. Le passage de l'une à l'autre coûte, en degré $i$, exactement
le facteur $(-1)^{i-1}(i-1)!$ --- et ce facteur, qui est le nombre annoncé plus
haut, traverse le dossier d'un bout à l'autre.

Il y apparaît sous quatre visages, sur des feuillets qui n'ont visiblement pas
été écrits ensemble. Comme la classe de Chern d'un point plongé dans un espace
projectif de dimension $r$, qui vaut $(-1)^{r-1}(r-1)!$ fois la classe de ce
point (page 9). Comme le défaut d'un certain isomorphisme, qui n'en est pas un,
entre deux façons de graduer l'anneau des classes (page 21). Comme la raison
pour laquelle une classe non nulle peut avoir une classe de Chern totale
parfaitement triviale --- il suffit d'une classe de torsion tuée par $(r-1)!$,
et c'est le contre-exemple de la page 13. Et enfin, quoique aucun feuillet ne
l'écrive, comme le coefficient qui relie le logarithme de la classe de Chern
totale au caractère de Chern, dont la page 43 calcule les deux premiers termes.
Quand Grothendieck, quelques années plus tard, cherchera un
« Riemann-Roch sans dénominateurs », c'est contre ce facteur qu'il se battra.

Le dossier lui-même n'est pas un texte. Il est un paquet : cinq chemises
inscrites de sa main le découpent en cinq courses --- « Classique pour les
Courbes », « Contre-exemples », « Cas des schémas », « Construction des Chern
covariants », « R.R. des surfaces » --- écrites sur trois sortes de papier,
dont les versos d'un texte dactylographié qui n'est pas de lui. Et ces courses
ne peuvent pas être contemporaines : certaines ne connaissent que les faisceaux
et les classes de Chern, tandis que d'autres emploient des mots --- topos,
complexe parfait, catégorie dérivée --- qui n'existaient pas encore lorsque le
théorème a été démontré. L'inventaire range le tout sous une date unique,
antérieure à 1958 ; les feuillets montrent que cela ne peut pas valoir pour
tous. Ce que l'on peut établir de ces écarts, et ce que l'on ne peut pas en
conclure, fait l'objet d'une section à part, immédiatement après celle-ci ;
c'est, pour un lecteur d'archives, la question centrale du dossier.

Il faut enfin dire ce que le dossier n'est pas. Ce ne sont pas des notes de
lecture, et ce n'est pas non plus une rédaction : c'est un homme qui recommence
le même calcul dans plusieurs langues successives, qui se pose des questions
par écrit --- « Est-ce que $\varphi_{K}$ est toujours bijectif ??? » --- et qui
y répond, deux lignes plus bas, « Non », en encadrant le mot. Plusieurs
feuillets s'arrêtent au milieu d'une phrase, un autre porte un titre sous lequel
rien n'est écrit, et le dernier s'interrompt six lignes après son début, à une
identification près de la formule d'adjonction qu'il allait obtenir. Le dossier
ne se termine pas : il cesse.

Les noms modernes sous lesquels chercher la suite : théorème de
Grothendieck-Riemann-Roch, caractère de Chern et classe de Todd, groupe de
Grothendieck d'un schéma, $\lambda$-anneaux et filtration $\gamma$,
Riemann-Roch sans dénominateurs, complexes parfaits et $K$-théorie à supports,
formule d'adjonction.

\keywords{Riemann-Roch theorem, Grothendieck-Riemann-Roch, Chern character, Todd class, total Chern class, Grothendieck group, K-theory of coherent sheaves, divisor class group, Picard group, Euler characteristic of a coherent sheaf, canonical class, lambda-ring, gamma-filtration, associated graded of the gamma-filtration, Chern ring, splitting principle, Riemann-Roch without denominators, regular closed immersion, normal bundle, projective space bundle, torsion in K-theory, ringed topos, equivariant module, induction product, symmetric group, tensor power operations, wreath product, projection formula, perfect complex, derived category, K-theory with supports, family of supports, Chow ring with supports, covariant Chern character, proper pushforward, higher direct image, fibration, degree of a covering, normalisation of a curve, adjunction formula, genus}
\end{resume}

\pagerange{1}{1}
\section*{Le fil du dossier, et les conventions}

\emph{Les stations.} Le dossier est un paquet de courses, séparées par cinq
chemises qu'il a inscrites lui-même (pages 1, 8, 19, 31 et 39) et que les
titres de section reprennent ici. Voici ce que chacune contient.

\begin{enumerate}
  \item \emph{Pages 2 et 3}, sous « Classique pour les Courbes » :
  Riemann-Roch sur une courbe non singulière, écrit dans l'anneau $E$ des
  classes de faisceaux cohérents. C'est la $K$-théorie avant la lettre $K$.
  \item \emph{Pages 4, 6 et 7} : le même énoncé refait en $K$-théorie, pour un
  schéma noethérien de dimension $\leqslant 1$, avec les deux variances
  $K^{\bullet}$ et $K_{\bullet}$, la filtration $\gamma$, $\operatorname{Pic}$
  et $\pi_{0}$.
  \item \emph{Pages 9, 11 et 13}, sous « Contre-exemples » : la classe de Chern
  totale d'une image directe le long de $X \to X \times \mathbb{P}^{r}$, le
  facteur $(r-1)!$, et une classe non nulle dont la classe de Chern totale
  vaut $1$.
  \item \emph{Pages 15 à 18} : les Modules sur un topos annelé où un groupe
  opère, le produit d'induction sur les groupes symétriques, et les formules
  pour les puissances tensorielles $T^{n}$.
  \item \emph{Pages 20 à 29}, sous « Cas des schémas » : le formalisme des
  $\lambda$-anneaux. L'anneau de Chern $CK$, le gradué $GK$ de la filtration
  $\gamma$, la comparaison entre les deux, le facteur $(-1)^{i-1}(i-1)!$, et un
  contre-exemple à la bijectivité.
  \item \emph{Pages 32 à 37}, sous « Construction des Chern covariants, et
  $K^{\bullet}$ et $A^{\bullet}$ à supports » : les catégories dérivées de
  complexes parfaits, la $K$-théorie et les groupes de Chow à supports, et
  Riemann-Roch conçu comme la mesure d'un défaut de covariance.
  \item \emph{Pages 40 à 44}, sous « R.R. des surfaces, et autres cas
  particuliers divers » : le carré de Grothendieck-Riemann-Roch, ses formules
  (1), (1 bis), (2) et (1 ter), puis les exemples --- une fibration, une courbe
  dans une courbe, un revêtement, une courbe dans une surface.
\end{enumerate}

\emph{Ce que seule une lecture d'ensemble peut dire.} Quatre liens, qu'aucun
feuillet n'énonce, et dont chacun enjambe au moins une chemise.

Le premier : le facteur $1 - \tfrac{1}{2}K$ par lequel la page 3 multiplie la
classe d'un faisceau \emph{est} la classe de Todd $\tau_{h}(X) = 1 -
\tfrac{1}{2}k$ de la page 42. Le même nombre, sur le même objet --- une courbe
---, à trente-neuf feuillets et au moins une campagne d'écriture de distance,
sous deux notations dont aucune ne renvoie à l'autre.

Le deuxième : l'opérateur $K_{1}$ de la page 3, défini comme la forme linéaire
nulle sur $\mathbb{Z}$ et égale au degré sur les classes de diviseurs, est
l'opérateur $K_{p}$ de la page 41 et le $K_{2}$ de la page 44. Un seul
opérateur : prendre la composante de degré $p$, puis son degré. Il est écrit
ici $\deg$ partout.

Le troisième, et c'est le plus lourd : le $(r-1)!$ des pages 9 à 13 et le
$(i-1)!$ de la page 21 sont le même facteur, et les deux contre-exemples qu'ils
servent à construire sont le même contre-exemple. Sur les feuillets
géométriques, une classe de diviseur de torsion tuée par $(r-1)!$ donne une
classe non nulle de classe de Chern totale triviale. Sur les feuillets
algébriques, l'anneau $CK$ a de la torsion là où le gradué $GK$ n'en a pas, et
la flèche de l'un vers l'autre n'est donc pas injective. C'est une seule et
même chose, dite deux fois et jamais rapprochée.

Le quatrième est de cette lecture et non du dossier. Les trois manifestations
précédentes sont des cas particuliers d'une identité que personne n'écrit ici :
si $c$ est la classe de Chern totale et $\operatorname{ch}$ le caractère de
Chern, alors
\[
  \log c \;=\; \sum_{m \geqslant 1} (-1)^{m-1}(m-1)!\,\operatorname{ch}_{m} .
\]
La page 43 en calcule, sans la nommer, les deux premiers termes, en écrivant la
classe de Chern totale comme l'exponentielle de
$\operatorname{ch}_{1} - \operatorname{ch}_{2}$.\footnote{L'identité se lit sur
les racines de Chern : $\log \prod_{j}(1 + a_{j}) = \sum_{m}
\frac{(-1)^{m-1}}{m}\sum_{j} a_{j}^{m}$, et $\sum_{j} a_{j}^{m} = m!\,
\operatorname{ch}_{m}$. Elle n'est écrite sur aucun feuillet du dossier ; elle
est donnée ici parce qu'elle est ce qui fait des trois observations
précédentes une seule.}

\emph{Deux liens plus faibles, et il faut dire pourquoi ils sont plus faibles.}
Les pages 15 à 18 construisent les opérations $\lambda^{i}$ par puissances
tensorielles à opérateurs, c'est-à-dire exactement les opérations dont les
pages 20 à 29 se donnent les axiomes ; mais aucune des deux courses ne renvoie
à l'autre, et l'ordre des feuillets n'est pas un argument. Et la filtration
$\gamma$ figure sur la page 4 comme sur la page 20 : elle ne peut donc pas
servir à distinguer les campagnes, contrairement à ce qu'on pourrait croire en
ne lisant que la seconde moitié du dossier.

\emph{Les conventions.} Sa notation change d'une course à l'autre ; une seule
est tenue ici, et les variantes sont signalées en note à l'endroit où elles
apparaissent.

\begin{itemize}
  \item Le caractère de Chern est $\operatorname{ch}$. Il écrit $c = \mathrm{ch}$
  aux pages 4 et 6, $\operatorname{ch}$ aux pages 36 et 37, et $\varphi$
  indexé par la base aux pages 40 à 44.
  \item La classe de Todd est $\operatorname{Todd}$. Il écrit $1 - \tfrac12 K$
  à la page 3 et $\tau_{h}$ aux pages 40 à 44.
  \item Les classes de Chern sont $c^{i}$, la classe totale $c$, et
  $\widetilde{c}(x) = [\varepsilon(x),\, 1 + c^{1}(x) + c^{2}(x) + \cdots]$
  lorsque le rang est porté à côté. Pour un anneau gradué $A$, on note
  $\widetilde{A} = \mathbb{Z} \times (1 + A^{\geqslant 1})$, muni de l'addition
  qui est la multiplication des séries et du produit $\ast$ décrit à la
  section sur les pages 20 à 29. C'est l'anneau où vivent les classes totales,
  et c'est le $\widetilde{A}$ des pages 9, 24, 27 et 40.
  \item $\deg$ est l'opérateur « composante de degré $p$, puis degré » :
  ses $K_{1}$, $K_{p}$, $K_{2}$.
  \item Les classes canoniques sont $k_{X}$, $k_{Y}$. Il écrit $K$ à la page 3,
  $k$ et $k'$ à la page 42, et $k_{x}$, $k_{y}$, $k_{X}$, $k_{Y}$ aux pages 42
  à 44.
  \item Aux pages 2 et 3, les fibrés vectoriels sont $V$, $V'$, $V''$ : la
  lettre $E$ y désigne déjà l'anneau des classes, et le feuillet l'emploie pour
  les deux.
  \item La lettre $E$ a trois emplois dans le dossier, dans trois courses
  distinctes : l'anneau des classes (pages 2 et 3), les éléments distingués
  $E_{s}$ d'un $\lambda$-anneau (pages 27 à 29), un fibré vectoriel (pages 41
  à 43). Aucun rapport entre les trois.
  \item La lettre $\xi$ en a trois aussi : $\xi = 1 - L^{-1}$ dans
  $K(\mathbb{P}^{r})$ (pages 9 à 13, conservé), le générateur de $CK$ à la page
  21 (écrit $\zeta$ ici), et la classe $(n,d) \in K(X)$ aux pages 42 et 43
  (conservé). La page 42 emploie de plus $\xi$ pour la classe du faisceau
  tordant, sur le même feuillet que l'emploi précédent : cette dernière est
  écrite $h$ ici.
\end{itemize}

\pagerange{1}{44}
\section*{Les campagnes du dossier : ce que les feuillets montrent}

Trois lecteurs indépendants, travaillant sur trois tranches différentes, ont
relevé la même chose : le dossier ne peut pas être homogène. Il faut le dire de
front, et il faut séparer soigneusement ce qui s'observe sur les feuillets, ce
qui s'en déduit, et ce qui n'en résulte pas.

\emph{Ce qui s'observe.} D'abord la matière. Le dossier est un paquet, non un
cahier : cinq chemises de sa main le découpent (pages 1, 8, 19, 31, 39), et les
courses qu'elles ouvrent sont écrites sur des supports différents --- feuillets
blancs libres, feuillets bleu-gris, et versos d'un texte dactylographié qui
n'est pas de lui, lequel occupe à lui seul huit des quarante-quatre pages.
Ensuite le vocabulaire, et c'est là que la chose se tranche.

Certaines pages emploient des mots dont la date de naissance est connue et
postérieure à 1958.

\begin{itemize}
  \item Page 15 : « topos annelé ». Le mot \emph{topos} est du séminaire de
  1963-1964, et n'a pas de sens avant lui.
  \item Page 32 : $D_{\mathrm{parf}}(f)$, « complexes de
  $\mathcal{O}_{X}$-Modules parfaits relativement à $f$ »,
  $K_{\mathrm{parf}}(f)$, ainsi que $\mathrm{R}i_{*}$ et les catégories
  $D_{Z}(X)$. Les catégories dérivées sont du début des années soixante ; les
  complexes parfaits relatifs à un morphisme sont l'objet propre du séminaire
  de 1966-1967 sur Riemann-Roch.
  \item Pages 32 à 36 : les familles de supports $\Phi$ en $K$-théorie,
  $K_{\Phi}(X)$, $A_{\Phi}(X)$, et les classes de Chern entre les deux.
  \item Pages 4 et 34 : le mot « préschéma », et le mot « Module » avec une
  capitale pour un faisceau de modules. Ce sont deux conventions des
  \emph{Éléments}, dont le premier volume est de 1960.
\end{itemize}

D'autres pages ne portent rien de tel. Les pages 2 et 3 écrivent $E$, et non
$K$, pour l'anneau des classes ; elles parlent de « faisceau algébrique
cohérent » et de « courbe non singulière », et leur seul outil est la classe de
Chern $c^{1}$ à valeurs dans le groupe des classes de diviseurs. Les pages 9 à
13 et 40 à 44 calculent avec la classe de Chern totale, le caractère de Chern,
la classe de Todd, $f_{!}$ et $f_{*}$, dans un anneau de Chow tensorisé par
$\mathbb{Q}$ : c'est l'appareil du théorème tel qu'il a été exposé en 1957 et
rédigé en 1958, et l'on n'y trouve ni complexe, ni support, ni catégorie
dérivée.

\emph{Ce qui s'en déduit.} Le dossier contient au moins deux campagnes, et la
datation « antérieures à 1958 » ne peut pas couvrir les deux. C'est une
conclusion sûre : elle ne repose pas sur une impression d'écriture mais sur des
mots dont on sait qu'ils n'existaient pas.

Un point mérite d'être ajouté, parce qu'il dérange la coupure la plus naturelle.
La frontière ne tombe pas sur les chemises. La première chemise, « Classique
pour les Courbes », contient à la fois les pages 2 et 3, qui ignorent la lettre
$K$, et les pages 4, 6 et 7, qui emploient $K^{\bullet}$ et $K_{\bullet}$ avec
leurs deux variances, $\pi_{0}(X)$, $\operatorname{Pic}(X)$, la filtration
$\gamma$ et le mot « préschéma ». Autrement dit, une couture passe à
l'intérieur d'une course que sa chemise donne pour une. Et la filtration
$\gamma$, dont on pourrait croire qu'elle marque le tard, est déjà sur la page
4.

\emph{Ce qui n'en résulte pas.} Trois choses, qu'il serait facile et faux de
croire acquises.

Aucun feuillet n'est daté par cette lecture. Que les pages 40 à 44 ne
contiennent aucune catégorie dérivée ne prouve nullement qu'elles soient
anciennes : un calcul classique s'écrit à n'importe quelle date, et ces pages
pourraient tout aussi bien être une reprise tardive, pour un cours ou pour une
rédaction. Ce que l'on peut dire est exactement ceci : les courses classiques
ne portent rien qui soit postérieur à 1958, et les courses tardives portent
quelque chose qui ne peut pas lui être antérieur. La première moitié de cette
phrase est beaucoup plus faible que la seconde, et c'est volontaire.

L'inventaire n'est pas pris en défaut. Sa date reste ici telle quelle,
$[1957]$, entre les crochets qui sont la convention des archivistes pour une
date déduite et non lue ; l'inférence appartient au lecteur, et cette édition
ne la fait pas. On notera seulement, comme un fait et non comme une thèse, que
les mots de l'inventaire --- « Notes Antiques (antérieures à 1958) » --- ne
figurent sur aucune des cinq chemises transcrites, qui portent respectivement
« Classique pour les Courbes », « Contre-exemples », « Cas des schémas »,
« Construction des Chern covariants » et « R.R. des surfaces ». D'où ils
viennent, la transcription ne le dit pas. Si l'étiquette est de lui, elle
désigne un paquet et non chacune de ses pièces, et un paquet ainsi étiqueté est
précisément ce à quoi l'on ajoute des feuillets plus tard --- mais ceci est une
conjecture, elle est signalée comme telle, et le dossier ne la tranche pas.

Enfin, et c'est ce qui importe le plus pour qui vient lire les mathématiques :
rien de ce qui suit ne dépend de la réponse. Les courses classiques démontrent
le théorème de 1957 ; les courses tardives construisent la machinerie qui
portera son nom dix ans après. La question qu'elles traitent, elle, est la même
de bout en bout --- la distance entre la classe de Chern et le caractère de
Chern, qui est le facteur $(-1)^{i-1}(i-1)!$ ---, et c'est vraisemblablement
pour cela que ces feuillets ont fini dans une même chemise. Qu'un feuillet ait
été écrit en 1957 ou en 1966 change l'histoire du texte ; cela ne change aucun
de ses énoncés.

\pagerange{2}{3}
\section*{Classique pour les Courbes : Riemann-Roch avant la lettre $K$ (pages 2 et 3)}

Soit $X$ une courbe algébrique non singulière, irréductible et
propre.\footnote{Le feuillet dit seulement « $X$ courbe non singulière ».
L'irréductibilité et la propreté sont indispensables : sans la première le
groupe $\pi_{0}$ intervient, sans la seconde la caractéristique
d'Euler-Poincaré de la formule (1) n'est pas finie. Elles sont ajoutées ici.}

\subsection*{L'anneau des classes}

On considère le groupe libre sur les faisceaux cohérents sur $X$, et l'on y
impose une relation par suite exacte courte : pour toute
$0 \to F' \to F \to F'' \to 0$, la classe de $F$ est la somme de celles de
$F'$ et de $F''$. Le produit tensoriel en fait un anneau, que le feuillet note
$E$. C'est le groupe que l'on appelle aujourd'hui $K_{0}(X)$ ; la lettre $K$ ne
figure pas sur ces deux pages.\footnote{La page 2 se sert de la lettre $E$ à la
fois pour l'anneau et pour un fibré vectoriel, dans la même proposition. Les
fibrés sont écrits $V$, $V'$, $V''$ ici.}

Le feuillet distingue dans $E$ les sous-objets suivants.

\begin{itemize}
  \item $N$, le sous-groupe engendré par les éléments $(F) - (F') - (F'')$ ---
  c'est-à-dire, dans le groupe libre, ce que l'on quotiente.\footnote{Le
  feuillet travaille donc simultanément dans le groupe libre et dans son
  quotient, sans le dire. Les énoncés $E/N \simeq \cdots$ n'ont de sens que
  dans le premier ; les lignes qui identifient $E$ à l'anneau des classes n'ont
  de sens que dans le second. La lecture suit cet usage, qui est cohérent
  pourvu qu'on prenne $E$ pour le groupe libre et $E/N$ pour l'anneau des
  classes proprement dit.}
  \item $\mathcal{D}$, le sous-groupe engendré par les faisceaux gratte-ciel
  $k_{x}$ : c'est le groupe des diviseurs.
  \item $\mathcal{L} \subset \mathcal{D}$, le groupe des diviseurs
  linéairement équivalents à $0$.
  \item $\Phi$, la sous-algèbre engendrée par les faisceaux localement libres,
  isomorphe à l'algèbre des classes de fibrés vectoriels, et $\Phi_{N}$ son
  intersection avec $N$.
  \item $E_{0}$, le noyau d'une augmentation dont le feuillet ne laisse pas
  lire le nom.\footnote{Le mot, abrégé, commence par un $g$ et se termine en
  \texttt{-ti.} ; il n'a pas été lu. Le rôle que $E_{0}$ joue dans les deux
  identités $\mathcal{D} + N = E_{0}$ et $E = \Phi + \mathcal{D}$ est celui du
  noyau du rang, mais la lecture n'est pas reconstruite ici.}
\end{itemize}

Le feuillet affirme alors, en marquant chaque ligne d'une étoile :
\[
  E = \Phi + \mathcal{D}, \qquad \Phi \cap \mathcal{D} = 0,
  \qquad E/N \;\simeq\; \Phi/\Phi_{N} \;\simeq\; \mathbb{Z} + \mathfrak{P} ,
\]
où $\mathfrak{P}$ est le groupe des classes de diviseurs, c'est-à-dire
$\mathcal{D}/\mathcal{L} = \operatorname{Pic}(X)$, et où l'isomorphisme envoie
$(1,0)$ sur la classe de $\mathcal{O}_{X}$ et induit sur $\mathfrak{P}$ le
passage au quotient de l'inclusion $\mathcal{D} \subset E$.\footnote{Sur le
feuillet, $\mathcal{L}$, $\mathcal{D}$ et $\mathfrak{P}$ ne se distinguent que
par le tracé de la lettre, et le transcripteur signale que ces lectures sont
solidaires : si $\mathfrak{P}$ n'est pas le groupe des classes de diviseurs,
elles bougent toutes. L'algèbre les fixe pourtant sans recours au tracé. La
page 3 pose $\mathcal{D} \cap N = \mathcal{L}$, donc
$E_{0}/N = (\mathcal{D} + N)/N \simeq \mathcal{D}/\mathcal{L}$ : le quotient est
forcément le groupe des classes de diviseurs, et c'est ce groupe qui reçoit
$c^{1}$ et qui contient la classe canonique $K$ de la formule (1). Aucune autre
lecture ne rend ces quatre emplois compatibles.}

\subsection*{La proposition, et ce qu'elle sert à établir}

\textbf{Proposition (page 2).} \emph{Le sous-groupe $N$ est engendré par les
éléments}
\[
  c(V) - c(V') - c(V''), \qquad
  c(V) - c(W) - c\bigl(\mathcal{O}_{X}(V)/f\,\mathcal{O}_{X}(W)\bigr),
\]
\emph{où la première famille court sur les suites exactes de fibrés
vectoriels et la seconde sur les couples de fibrés en droites munis d'une
section} $f \in H^{0}(X, \mathcal{O}_{X}(W^{*} \otimes V))$ \emph{non
nulle.}\footnote{Le feuillet laisse illisible le mot qualifiant $f$ ; le membre
de droite n'a de sens que si $f$ est injective, ce qu'une section non nulle
assure sur une courbe intègre. La restriction aux fibrés de rang $1$ est celle
du feuillet.}

Le feuillet reformule aussitôt : cela revient à dire que $N$ est engendré par
les $\lambda(D + D') - \lambda(D) - \lambda(D')$, où $\lambda(\Delta)$ est la
classe du faisceau inversible associé au diviseur $\Delta$, et qu'on peut s'y
ramener au cas $D \geqslant 0$.

La page 3 en tire quatre identités équivalentes :
\[
  \mathcal{D} + N = E_{0}, \qquad \mathcal{D} \cap N = \mathcal{L},
  \qquad \Phi \cap N = \Phi_{N} ,
\]
la quatrième, $\Phi + N = E$, étant biffée sur le feuillet.\footnote{Elle est
d'ailleurs fausse telle quelle : $\Phi + N$ ne contient pas les classes de
faisceaux de torsion, qui engendrent $\mathcal{D}$. La raturer est la bonne
décision, et la lecture la respecte.} La deuxième est le cœur de l'affaire :
elle dit que deux diviseurs dont les classes coïncident dans l'anneau des
classes sont linéairement équivalents, et c'est exactement ce que la seconde
famille de générateurs de la proposition fournit, puisque celle-ci n'identifie
deux diviseurs que lorsqu'une fonction les relie.\footnote{Le feuillet écrit
les éléments de $\mathcal{D} \cap N$ « diviseurs « $f$-équivalents
» à $0$ », le préfixe entre guillemets étant une haste bouclée
que le transcripteur ne tranche pas entre $f$ et $l$ (pour
« linéairement »). Le contenu de la proposition rend les deux lectures
équivalentes : $f$-équivalent, c'est-à-dire équivalent au moyen d'une fonction,
\emph{est} linéairement équivalent.}

\subsection*{Riemann-Roch en une ligne}

On définit $E \to \mathbb{Z} + \mathfrak{P}$ ainsi : sur $\mathcal{D}$, c'est
l'application canonique $\mathcal{D} \to \mathfrak{P}$ ; sur $\Phi$, un fibré
vectoriel $V$ de rang $n$ va sur $(n,\, c^{1}(V))$. C'est, dans le langage
d'aujourd'hui, le caractère de Chern : pour une courbe,
$\operatorname{ch} = (\operatorname{rg}, c^{1})$, et il est à valeurs entières
puisque les termes de degré $\geqslant 2$ n'existent pas.

La formule de la page 3, pour $F$ un faisceau cohérent, est alors
\begin{equation*}
  \chi(F) \;=\; \deg\Bigl( \operatorname{ch}(F)\cdot
  \bigl(1 - \tfrac{1}{2}k_{X}\bigr) \Bigr)
  \tag{1}
\end{equation*}
où $k_{X} \in \mathfrak{P}$ est la classe canonique, et où $\deg$ est la forme
linéaire nulle sur $\mathbb{Z}$ et égale au degré sur
$\mathfrak{P}$.\footnote{Il écrit le membre de gauche avec un $s$ long bouclé,
que la transcription rend par $\int$ : c'est la caractéristique
d'Euler-Poincaré, et elle est écrite $\chi$ ici. Il écrit $K$ pour la classe
canonique et $K_{1}$ pour la forme linéaire ; ces deux lettres sont remplacées,
la première parce que $K$ désigne partout ailleurs dans le dossier un
$\lambda$-anneau ou un groupe de Grothendieck, la seconde parce que le même
opérateur reparaît aux pages 41 et 44 sous le nom $K_{p}$.}

C'est le théorème de Riemann-Roch. Il suffit en effet, comme le feuillet le
note, de le vérifier sur $F = k_{x}$ et sur $F = \mathcal{O}_{X}$, tout
faisceau cohérent sur une courbe ayant même classe qu'une combinaison de
ceux-là. Pour $\mathcal{O}_{X}$ : $\operatorname{ch} = (1,0)$, le produit vaut
$(1, -\tfrac12 k_{X})$, et son degré est $-\tfrac12(2g-2) = 1-g = \chi(\mathcal{O}_{X})$.
Pour $k_{x}$ : $\operatorname{ch} = (0, [x])$, le produit vaut $(0,[x])$, de
degré $1 = \chi(k_{x})$. Pour $F$ quelconque de rang $n$ et de degré $d$, on
trouve $d + n(1-g)$.

Le facteur $1 - \tfrac12 k_{X}$ est la classe de Todd de la courbe. Il
reviendra page 42, à l'autre bout du dossier, calculé cette fois comme
$-k/(1 - e^{k})$ et noté $\tau_{h}(X)$ : c'est le même nombre, et rien sur les
feuillets ne les rapproche.

\pagerange{4}{7}
\section*{Le même énoncé en $K$-théorie : dimension $\leqslant 1$ (pages 4, 6 et 7)}

Les trois feuillets suivants refont la page 3 dans un autre vocabulaire. Soit
$X$ un schéma noethérien de dimension $\leqslant 1$ admettant un faisceau
inversible ample.\footnote{Le feuillet écrit « préschéma », qui est la
convention des \emph{Éléments} ; et il commence par écrire « admettant aussi un
Module inversible ample », le mot « aussi » étant biffé. L'hypothèse d'amplitude
n'est pas décorative : c'est elle qui donne le théorème de structure des fibrés
utilisé plus bas.} On dispose de deux groupes de Grothendieck, celui des fibrés
vectoriels $K^{\bullet}(X)$, contravariant, et celui des faisceaux cohérents
$K_{\bullet}(X)$, covariant.\footnote{Il les écrit avec un point surélevé et un
point souscrit. Ce sont les groupes que l'on note aujourd'hui $K^{0}(X)$ et
$K_{0}(X)$.}

\subsection*{Le côté contravariant}

$K^{\bullet}(X)$ est augmenté vers $H^{0}(X, \mathbb{Z}) = \mathbb{Z}^{\pi_{0}(X)}$
par le rang, qui est localement constant. Soit $I(X)$ l'idéal d'augmentation.
Pour $x \in I(X)$ on a $\gamma^{i}(x) = 0$ dès que $i \geqslant 2$, d'où
$F^{i}_{\gamma}(K^{\bullet}(X)) = 0$ pour $i \geqslant 2$, et par
conséquent
\[
  K^{\bullet}(X) \;\simeq\; \mathbb{Z}^{\pi_{0}(X)} \times \operatorname{Pic}(X),
  \qquad
  \operatorname{Pic}(X) \simeq \prod_{i \in \pi_{0}(X)} \operatorname{Pic}(X_{i}),
\]
la classe d'un fibré $V$ étant $\bigl[\operatorname{rg}(V),\, \det V\bigr]$,
c'est-à-dire
\[
  \operatorname{cl}(V) \;=\; \sum_{i} \bigl(\operatorname{rg}_{i}(V) - 1\bigr)\,1_{i}
  \;+\; \operatorname{cl}(\det V).
\]
La raison est que sur un tel $X$ tout fibré vectoriel est somme directe d'un
fibré trivial et de son déterminant.\footnote{Le feuillet ne donne pas cette
raison et laisse illisible la phrase qui suivait « faisceau loc. libre de rang
», ainsi que le rang lui-même, lu $0$. C'est le théorème de structure des
fibrés sur un schéma de dimension $\leqslant 1$ à faisceau ample, et c'est là
que sert l'hypothèse d'amplitude.} On pose $B^{\bullet}(X) = \mathbb{Z}^{\pi_{0}(X)}
\times \operatorname{Pic}(X)$ ; avec ces identifications, le caractère de
Chern $K^{\bullet}(X) \to B^{\bullet}(X)$ est l'identité.

\subsection*{Le côté covariant}

$K_{\bullet}(X)$ est augmenté vers $\mathbb{Z}^{\pi_{0}^{(1)}(X)}$, où
$\pi_{0}^{(1)}(X)$ est l'ensemble des composantes irréductibles de dimension
$1$, par les longueurs aux points génériques. Le noyau est engendré par les
classes de faisceaux à support fini et par les différences de faisceaux
inversibles sur les composantes de dimension $1$.\footnote{Le feuillet écrit ce
noyau $K_{0}(X) \simeq \mathbb{Z}^{k^{(0)}(X)} \times \operatorname{Pic}'(X)$
et glose le second facteur comme un produit de $K_{0}(X_{i})$ sur les
composantes de dimension $1$. Les indices ne sont pas tous lisibles et la
description exacte du premier facteur est perdue ; ce qui est donné ici est la
part de l'énoncé que le feuillet soutient.}

\subsection*{La flèche d'un côté à l'autre}

Les pages 6 et 7 sont un palimpseste rapide dont la prose de liaison est en
très grande partie perdue, et la seconde est barrée d'obliques depuis son
milieu.\footnote{Aucune tentative n'est faite ici de combler ces lacunes. Ce
qui suit est ce que portent les formules, non ce que disait le texte qui les
reliait.} Ce qu'elles établissent malgré tout est net.

En dimension $0$, le rang et la longueur sont des isomorphismes de
$\lambda$-anneaux $K^{\bullet}(X) \simeq \mathbb{Z}$ et $K_{\bullet}(X) \simeq
\mathbb{Z}$, et le caractère de Chern est l'identité. En dimension $1$, la page
6 dresse la liste des invariants : l'ensemble fini $\pi_{0}(X)$ des composantes
irréductibles, le groupe $\operatorname{Pic}(X)$, le groupe $K_{(0)} =
\prod_{i} K_{(0)}(X_{i})$, les entiers $\mu_{i} = \operatorname{lg}_{i}(\mathcal{O}_{X})$,
un élément $\delta \in K_{0}(X)$, et enfin des applications
$u_{i} : \operatorname{Pic}(X) \to K_{(0)}$, $L \mapsto
\operatorname{cl}\bigl((L-1)\mathcal{O}_{X_{i}}\bigr)$.

La page 7 écrit alors l'homomorphisme $K^{\bullet}(X) \to K_{\bullet}(X)$. Il
est déterminé par l'image de $1$, soit
$1_{\bullet} = \sum_{i} \mu_{i} 1_{i} + \delta_{\varepsilon}$, et vaut
\[
  u(n, \xi) \;=\; n \sum_{i} \mu_{i}\,1_{i}
  \;+\; \Bigl( n\,\delta_{\varepsilon} + \sum_{i} \mu_{i}\,u_{i}(\xi) \Bigr),
  \qquad n \in \mathbb{Z},\ \xi \in \operatorname{Pic}(X).
\]
Ce que les feuillets ne disent pas, et qui explique tout, est que cette flèche
est simplement la multiplication par la classe de $\mathcal{O}_{X}$ dans
$K_{\bullet}(X)$ : $K_{\bullet}$ est un module sur $K^{\bullet}$, et l'unique
morphisme de modules $K^{\bullet} \to K_{\bullet}$ est $x \mapsto
x \cdot [\mathcal{O}_{X}]$. La formule s'obtient alors en développant
$(n + \xi)\cdot 1_{\bullet}$ et en observant que $\xi \cdot \delta_{\varepsilon} = 0$
pour raison de degré, tandis que $\xi \cdot 1_{i} = u_{i}(\xi)$ par définition
de $u_{i}$.\footnote{C'est donc, en langage moderne, la flèche
$K^{0}(X) \to K_{0}(X)$, dont on sait qu'elle est bijective lorsque $X$ est
régulier. Le feuillet
demande, en dimension $0$, « pour $2$ choix ! », et l'annotation marginale
signale que l'invariant en cause dépend du choix de la classe de $1$ dans
$K_{\bullet}(X)$ : c'est exactement l'observation ci-dessus.}

\pagerange{9}{13}
\section*{Contre-exemples : la classe de Chern d'une immersion, et le facteur $(r-1)!$ (pages 9 à 13)}

Cette course pose une question précise : les classes de Chern d'une image
directe sont-elles calculables à partir de celles de la classe qu'on pousse ?
La réponse est oui, avec un facteur entier qui grossit ; et le prix de ce
facteur est qu'une classe non nulle peut avoir une classe de Chern totale
triviale. C'est là le contre-exemple qui donne son titre à la chemise.

\subsection*{La situation}

Soit $X$ un schéma, $r \geqslant 1$, et $i : X \to X \times \mathbb{P}^{r}$
l'immersion fermée d'une section constante. C'est une immersion régulière de
codimension $r$ à fibré normal trivial. On dispose de
\[
  K(X \times \mathbb{P}^{r}) \;=\;
  K(X) \otimes K(\mathbb{P}^{r}), \qquad
  K(\mathbb{P}^{r}) = \mathbb{Z}[\xi]/(\xi^{r+1}),\quad \xi = 1 - L^{-1},
\]
où $L = \mathcal{O}(1)$ ; la classe $\xi^{r}$ est celle du faisceau structural
d'un point.\footnote{Le feuillet écrit la présentation sous la forme
$K(X) \otimes K(\mathbb{P}^{r})/(\xi = 1 - L^{-1})$, c'est-à-dire en nommant le
générateur plutôt qu'en donnant la relation $\xi^{r+1} = 0$. Les deux disent la
même chose.} On a donc $i_{!}(x) = x \otimes \xi^{r}$, et la flèche
$K(X) \to K(X \times \mathbb{P}^{r})$ est injective, la composée avec
$p_{*}$ valant l'identité :

\begin{tikzcd}[column sep=large, row sep=small]
K(X) \arrow[r, "\text{inj.}"] \arrow[d, "\widetilde{c}_{X}"'] & K(X \times \mathbb{P}^{r}) \arrow[d, "\widetilde{c}_{X \times \mathbb{P}^{r}}"] \\
\widetilde{A}(X) \arrow[r] & \widetilde{A}(X \times \mathbb{P}^{r})
\end{tikzcd}

\subsection*{La classe de Chern totale d'un point}

Puisque $\xi^{r} = \sum_{i=0}^{r} \binom{r}{i}(-1)^{i} L^{-i}$ et
$c(L^{-i}) = 1 - i\eta$, avec $\eta$ la classe d'un hyperplan, on a
\[
  c(\xi^{r}) \;=\; \prod_{i=0}^{r} (1 - i\eta)^{(-1)^{i}\binom{r}{i}} .
\]
En passant au logarithme, le coefficient de $\eta^{m}$ est
$-\frac{1}{m}\sum_{i}(-1)^{i}\binom{r}{i} i^{m}$, somme qui est nulle pour
$m < r$ et vaut $(-1)^{r} r!$ pour $m = r$. Comme $\eta^{r+1} = 0$, il ne reste
qu'un terme et l'exponentielle se réduit à sa partie linéaire :
\[
  c(\xi^{r}) \;=\; 1 + (-1)^{r-1}(r-1)!\ \eta^{r} .
\]
C'est le premier visage du nombre du dossier : la classe de Chern de dimension
maximale du faisceau structural d'un point de $\mathbb{P}^{r}$ n'est pas la
classe de ce point, mais son multiple par
$(-1)^{r-1}(r-1)!$.\footnote{L'exposant de $(-1)$ est écrit $r+1$ à la page 9 et
$r-1$ à la page 11, et la transcription signale que la page ne tranche pas
entre les deux. Elle n'a pas à trancher : $(-1)^{r+1} = (-1)^{r-1}$. Les deux
leçons sont la même, et la variance relevée n'en est pas une.}

\subsection*{Les polynômes $Q^{(r)}_{i}$}

Le feuillet écrit ensuite $c(i_{!}(x)) = c(x) \ast c(\xi^{r})$ et développe. Le
résultat auquel il arrive est
\[
  c\bigl(i_{!}(x)\bigr) \;=\; 1 + (-1)^{r-1}(r-1)!\ \eta^{r}
  \Bigl[\, \varepsilon(x) + \sum_{i \geqslant 1}
  Q^{(r)}_{i}\bigl(c^{1}(x), \ldots, c^{i}(x)\bigr) \Bigr],
\]
où $Q^{(r)}_{i-r}$ est, par définition, la partie linéaire en la variable de
degré $r$ du polynôme universel $Q_{i}$ qui exprime $c^{i}$ d'un produit ---
les termes de degré $\geqslant 2$ en cette variable disparaissant puisque
$(\eta^{r})^{2} = 0$.\footnote{La page 9 développe $c(x) \ast (1 - i\eta)$ en
$\sum_{j \leqslant k} c^{j}(x)(-i\eta)^{k-j}$, sans les coefficients
binomiaux : le terme de degré $k$ de $c(x \otimes M)$, pour $M$ inversible de
première classe $m$, est $\sum_{j=0}^{k}\binom{n-j}{k-j} c^{j}(x) m^{k-j}$ avec
$n = \varepsilon(x)$. Les binômes sont rétablis ici. Le feuillet porte sous
cette ligne une accolade et une glose dont le milieu est illisible.}

Le feuillet ne calcule le coefficient dominant de $Q^{(r)}_{i}$ que pour
$i = 1$, où il trouve $-r\,c^{1}$. Sa propre méthode le donne pour tout $i$ :
\[
  Q^{(r)}_{i}\bigl(c^{1}, \ldots, c^{i}\bigr)
  \;=\; -\,\frac{(r+i-1)!}{(r-1)!\,(i-1)!}\ c^{i}
  \;+\; (\text{termes décomposables}),
\]
et la lettre du feuillet le confirme, puisqu'il écrit son coefficient sous la
forme $-\frac{(r+1-1)!}{(r-1)!\,(1-1)!}$, c'est-à-dire précisément cette
expression pour $i = 1$.\footnote{La formule générale est fournie ici et ne
figure pas sur le feuillet, qui n'écrit que le cas $i=1$. Voici d'où elle
vient. Le fibré normal étant trivial, $\operatorname{ch}(i_{!}x) = \eta^{r}
\operatorname{ch}(x)$, d'où $\operatorname{ch}_{r+j}(i_{!}x) = \eta^{r}
\operatorname{ch}_{j}(x)$ et, par l'identité $\log c = \sum_{m}(-1)^{m-1}(m-1)!
\operatorname{ch}_{m}$ et $(\eta^{r})^{2} = 0$,
$c(i_{!}x) = 1 + (-1)^{r-1}\eta^{r}\sum_{j}(-1)^{j}(r+j-1)!\operatorname{ch}_{j}(x)$.
On conclut par $\operatorname{ch}_{i} = \frac{(-1)^{i-1}}{(i-1)!}c^{i} +
(\text{décomposables})$, conséquence des identités de Newton. Le feuillet note
au passage $Q^{(r)}_{0} = 0$, ce qui est cohérent : le terme de rang est porté
à part, sous la forme du $\varepsilon(x)$ entre crochets.} Le feuillet ajoute,
en marge de ce calcul, l'avertissement que les coefficients numériques
proviennent tous du $(r-1)!$ : c'est exact, et c'est ce qui rend le paragraphe
suivant possible.

Il note enfin cette situation $c(i_{!}(x)) = (i_{*})_{r}\bigl(c(x)\bigr)$, en
introduisant donc un « $i_{*}$ » agissant directement sur les classes de Chern
totales. C'est la question que l'on appellera plus tard Riemann-Roch sans
dénominateurs : obtenir une formule d'image directe pour les classes de Chern
entières, et non pour le caractère de Chern rationnel.

\subsection*{Le contre-exemple}

Soit $Y \subset X$ un diviseur, de classe $x' \in A^{1}(X)$. On a
$c(\mathcal{O}_{Y}) = (1 - x')^{-1} = \sum_{n} (x')^{n}$. En appliquant ce qui
précède à $x = [\mathcal{O}_{Y}]$, chaque terme du crochet porte en facteur
$(r-1)!$ multiplié par une puissance de $x'$.\footnote{Explicitement, le terme
de rang $j$ vaut $(r+j-1)!\operatorname{ch}_{j}(\mathcal{O}_{Y})/(r-1)! =
\binom{r+j-1}{j}(-1)^{j-1}(x')^{j}$ au signe près, dont tous les coefficients
sont des entiers multipliés par $(r-1)!$ une fois le facteur global rétabli.}
Par conséquent :

\textbf{Si $(r-1)!\,x' = 0$ dans $A^{1}(X)$, alors
$c\bigl(i_{!}(\mathcal{O}_{Y})\bigr) = 1$.}\footnote{Le dernier membre de la
page 13 est lu $0$ ; la classe de Chern \emph{totale} vaut $1$, et c'est sa
partie de degré $\geqslant 1$ qui vaut $0$. Les deux énoncés disent la même
chose et ne diffèrent que par ce que $\widetilde{C}$ désigne ; le feuillet ne
tranche pas, et il n'y a rien à trancher.}

Autrement dit, la classe non nulle $[\mathcal{O}_{Y \times \mathbb{P}^{r}}]$ a
une classe de Chern totale triviale : les classes de Chern ne séparent pas.
Reste à produire un $x' \neq 0$ tué par $(r-1)!$, ce qu'une classe de torsion
de $\operatorname{Pic}^{0}$ fournit --- par exemple un point de torsion d'ordre
divisant $(r-1)!$ sur une courbe elliptique.\footnote{Le feuillet demande
seulement $(r-1)!\,x' = 0$, et la fin de sa dernière phrase est illisible.
L'en-tête de la transcription du lot nomme une courbe elliptique ; le corps
transcrit du feuillet ne la nomme pas. L'exemple est donc donné ici comme ce
qu'exige la condition, non comme ce que la page écrit.}

C'est le même phénomène que celui de la page 21, qui l'énoncera en pure algèbre :
$CK$ a de la torsion, $GK$ n'en a pas, et la flèche de l'un vers l'autre n'est
donc pas injective. Ici la torsion est celle de $\operatorname{Pic}^{0}$ ; là,
celle qu'introduisent les factorielles.

\pagerange{15}{18}
\section*{Modules sur un topos annelé où un groupe opère (pages 15 à 18)}

Ces quatre feuillets ne sont précédés d'aucune chemise et leur sujet n'est pas
celui de la course qui les précède.\footnote{Le titre de cette section est de
l'éditeur ; les mots sont ceux de la première ligne du feuillet 15. C'est la
seule course du dossier qui ne soit annoncée par aucune chemise inscrite.}
C'est, en apparence, une digression ; c'est en réalité la construction des
opérations dont la course suivante se donnera les axiomes, et c'est aussi la
course la plus tardive du dossier, puisqu'elle commence par le mot
\emph{topos}.

Soit $\mathcal{X}$ un topos annelé, d'anneau commutatif $\mathcal{O}$. Pour
tout groupe $G$ on note $A(\mathcal{X}, G)$ la catégorie des $\mathcal{O}$-Modules
de $\mathcal{X}$ munis d'une action de $G$ --- autrement dit, la catégorie des
$\mathcal{O}[G]$-Modules.

\subsection*{Restriction, induction, et le produit extérieur}

Un homomorphisme $u : G \to G'$ donne une restriction
$u^{*} : A(\mathcal{X}, G') \to A(\mathcal{X}, G)$ et une induction
$u_{*}$ en sens inverse. Lorsque $u$ est injectif d'indice fini, on a
l'isomorphisme canonique de $\mathcal{O}$-Modules
\[
  u^{*} u_{*}(F) \;\simeq\; \bigsqcup_{G'/G} F ,
\]
la somme portant sur les classes à gauche, $G$ permutant les facteurs par
translation.\footnote{Le feuillet écrit $F^{(G'/G)}$, notation qui se lit aussi
bien comme un produit. Pour un indice fini les deux coïncident ; l'hypothèse
d'indice fini n'est pas sur le feuillet et est ajoutée ici.} Les
sous-catégories pleines des Modules localement libres, libres de type fini, et
localement libres de type fini sur $\mathcal{O}[G]$ sont stables par ces
foncteurs.

On dispose dans $A(\mathcal{X}, G)$ du produit tensoriel, commutatif et
associatif, et des isomorphismes fonctoriels
\[
  u^{*}(F' \otimes G') \simeq u^{*}(F') \otimes u^{*}(G'),
  \qquad
  u_{*}\bigl(F \otimes u^{*}(G')\bigr) \simeq u_{*}(F) \otimes G' ,
\]
la seconde étant la formule de projection. Pour $F \in A(\mathcal{X}, G)$ et
$H$ un second groupe, $G \in A(\mathcal{X}, H)$, on pose
\[
  F \boxtimes G \;=\; \mathrm{pr}_{1}^{*}(F) \otimes \mathrm{pr}_{2}^{*}(G)
  \;\in\; A(\mathcal{X}, G \times H),
\]
et le produit tensoriel ordinaire se retrouve par restriction le long de la
diagonale : $F \otimes G = (\mathrm{diag}_{G})^{*}(F \boxtimes
G)$.\footnote{Le feuillet écrit l'objet dans $A(\mathcal{X}, H \times G)$ alors
que la définition le place dans $A(\mathcal{X}, G \times H)$ ; c'est une
inversion sans conséquence, les deux catégories étant canoniquement
isomorphes.}

\subsection*{Le produit d'induction, et les puissances tensorielles}

Soit $u_{pq} : \mathfrak{S}_{p} \times \mathfrak{S}_{q} \to
\mathfrak{S}_{p+q}$ l'injection canonique. Pour $F$ dans
$A(\mathcal{X}, \mathfrak{S}_{p})$ et $G$ dans $A(\mathcal{X}, \mathfrak{S}_{q})$
on pose
\[
  F \ast G \;=\; u_{pq\,*}\bigl(F \boxtimes G\bigr)
  \;\in\; A(\mathcal{X}, \mathfrak{S}_{p+q}),
\]
opération commutative et associative : elle fait de la somme des
$A(\mathcal{X}, \mathfrak{S}_{n})$ un anneau gradué.\footnote{Il abrège
$A(\mathcal{X}, p)$ pour $A(\mathcal{X}, \mathfrak{S}_{p})$. C'est le produit
d'induction, celui qui fait de la somme des catégories de représentations des
groupes symétriques l'anneau des fonctions symétriques.}

Notant $T^{n}(F)$ la puissance tensorielle $n$-ième munie de son action de
$\mathfrak{S}_{n}$, la page 18 énonce les trois formules
\[
  T^{n}(F + G) \simeq \bigsqcup_{p+q=n} T^{p}(F) \ast T^{q}(G),
  \qquad
  T^{m}\bigl(T^{n}(F)\bigr) \simeq w_{m,n}^{*}\bigl(T^{mn}(F)\bigr),
  \qquad
  T^{m}(F \otimes G) \simeq T^{m}(F) \otimes T^{m}(G),
\]
où $w_{m,n} : \mathfrak{S}_{n} \wr \mathfrak{S}_{m} \to \mathfrak{S}_{mn}$ est
le plongement du produit en couronne.\footnote{Le feuillet écrit $w_{m,n} :
\mathfrak{S}_{m} \to \mathfrak{S}_{?}$, la cible étant illisible, et le
qualifie de « convenable ». La source ne peut pas être $\mathfrak{S}_{m}$ :
$T^{m}(T^{n}F)$ est $T^{mn}(F)$ muni de l'action du groupe des permutations de
$mn$ objets rangés en $m$ paquets de $n$, qui est le produit en couronne. La
correction est faite ici.} La première est la formule du binôme pour les
puissances tensorielles, et c'est elle qui produit les opérations
$\lambda^{i}$ : le facteur $T^{p}(F) \ast T^{q}(G)$ est exactement ce qui
rend $\lambda$ additif au sens des $\lambda$-anneaux.

Le feuillet termine en isolant la condition ($\ast$) que doit remplir un Module
$F$ pour que les $u_{pq\,*}$ de ses composantes soient localement libres sur
les $\mathcal{O}[\mathfrak{S}_{p_{1}} \times \cdots \times
\mathfrak{S}_{p_{r}}]$, pour toute décomposition $\sum p_{i} = n$. C'est la
condition d'admissibilité qui fait fonctionner tout ce qui précède ; son
énoncé n'est pas lisible au-delà de sa forme.\footnote{La prose de cette
dernière ligne est presque entièrement perdue. Rien n'en est reconstruit ici.}

\pagerange{20}{29}
\section*{Cas des schémas : les $\lambda$-anneaux, $CK$ et $GK$ (pages 20 à 29)}

La chemise qui ouvre cette course porte, de sa main, le programme : « Filtration
des $\lambda$-anneaux, classes de Chern universelles ; $\lambda$-anneaux
« spéciaux» i.e. avec éléments
« géométriques» ; Application ». Le dossier
exécute les deux premiers points et pas le troisième.

\subsection*{Le cadre}

Soit $K$ un $\lambda$-anneau spécial, augmenté vers $\mathbb{Z}$ par
$\varepsilon$, muni d'une famille $(E_{s})_{s \in S}$ d'éléments
\emph{géométriques}, c'est-à-dire tels que
\[
  \varepsilon(E_{s}) = n_{s} \geqslant 0, \qquad
  \lambda^{i}(E_{s}) = 0 \ \text{ pour } i > n_{s}, \qquad
  \lambda^{n_{s}}(E_{s}) \ \text{ inversible.}
\]
Ce sont les axiomes que vérifie la classe d'un fibré vectoriel de rang $n_{s}$,
dont le déterminant est inversible ; d'où le nom. On demande de plus un
homomorphisme déterminant $\det : K \to K^{*}$ avec $\det(E_{s}) =
\lambda^{n_{s}}(E_{s})$.\footnote{Les pages 27 et 28 posent ces axiomes ; la
page 20 les annonce et la page 21 en donne, sous forme de proposition barrée
d'une grande croix, une variante équivalente en termes des $c^{i}$ : une
famille engendre $K$ pour la structure de $\lambda$-anneau avec déterminants
inversibles si et seulement si $c^{i}(E_{s}) = 0$ pour $i > \varepsilon(E_{s})$.
Il l'avait notée « à prouver directement ? » avant de la rayer.}

On note $I(K) = \ker \varepsilon$, $\gamma^{i}$ les opérations déduites des
$\lambda^{i}$ par $\gamma_{t} = \lambda_{t/(1-t)}$, et $K^{(i)}$ la filtration
$\gamma$ : le sous-groupe engendré par les produits
$\gamma^{i_{1}}(x_{1})\cdots\gamma^{i_{m}}(x_{m})$ avec $x_{j} \in I(K)$ et
$\sum i_{j} \geqslant i$. Le gradué associé est $GK = \operatorname{gr} K$, et
les classes de Chern sont
\[
  c^{i}(x) \;=\; \gamma^{i}\bigl(x - \varepsilon(x)\bigr) \bmod K^{(i+1)}
  \;\in\; GK^{i} .
\]

\subsection*{L'anneau $\widetilde{A}$ et le produit $\ast$}

Pour un anneau gradué $A$, on note $\widetilde{A} = \mathbb{Z} \times
(1 + A^{\geqslant 1})$, dont les éléments s'écrivent $[\varepsilon,\,
1 + a^{1} + a^{2} + \cdots]$. L'addition de $\widetilde{A}$ est l'addition des
rangs et la multiplication des séries ; sa multiplication, notée $\ast$, est
donnée par les polynômes universels $Q_{n}$ :
\[
  c^{n}(x \ast y) \;=\; Q_{n}\bigl(c^{1}(x), \ldots, c^{n}(x);\,
  c^{1}(y), \ldots, c^{n}(y)\bigr),
\]
lesquels sont déterminés par le principe de scindage : si $c(x) =
\prod_{i}(1 + a_{i})$ sur $\varepsilon(x)$ racines et $c(y) =
\prod_{j}(1 + b_{j})$ sur $\varepsilon(y)$ racines, alors $c(x \ast y) =
\prod_{i,j}(1 + a_{i} + b_{j})$. C'est l'anneau où vit la classe de Chern
totale, et c'est celui que le dossier note $\widetilde{A}$ aux pages 9, 24, 27
et 40.

\subsection*{L'anneau de Chern $CK$ et la comparaison avec $GK$}

La page 28 définit $CK$ par générateurs et relations : des générateurs
$c^{i}(x)$ pour $x \in K$ et $i \geqslant 1$, et les relations
\[
  \begin{cases}
    c^{i}(1) = 0 & i \geqslant 1, \\[3pt]
    c^{n}(x+y) = \displaystyle\sum_{i+j=n} c^{i}(x)\,c^{j}(y)
      & x, y \in K,\ n \geqslant 1, \\[8pt]
    c^{n}(x \ast y) = Q_{n}\bigl(c^{1}(x), \ldots ;\, c^{1}(y), \ldots\bigr)
      & x, y \in I(K),\ n \geqslant 1, \\[5pt]
    c^{n}(E_{s}) = 0 & n > n_{s}.
  \end{cases}
\]
La dernière relation est celle qui fait intervenir la donnée géométrique.
Puisque $\lambda^{i}(E_{s}) = 0$ pour $i > n_{s}$ entraîne $\gamma^{n}(E_{s} -
n_{s}) = 0$ pour $n > n_{s}$, les classes $\gamma^{i}(x - \varepsilon(x))$
satisfont à toutes ces relations, d'où un homomorphisme
\[
  \varphi_{K} : CK \longrightarrow GK, \qquad
  c^{i}(x) \longmapsto \gamma^{i}\bigl(x - \varepsilon(x)\bigr) \bmod K^{(i+1)},
\]
visiblement surjectif. Le feuillet note en passant --- et c'est l'observation
juste --- que plus la famille $(E_{s})$ est grande, plus $CK$ est petit, et plus
il y a de chances que $\varphi_{K}$ soit un isomorphisme.\footnote{Le bas de la
page 28 porte une longue note d'une plume plus fine dont seul le premier mot se
laisse lire, et une note marginale interrogeant l'itération du produit $\ast$
sur les $P(\gamma^{i}(x - \varepsilon(x)))$. Ni l'une ni l'autre n'est
reconstruite ici.}

La page 29 ajoute la réciproque du point de vue : tout homomorphisme de Chern
$\chi : K \to \widetilde{A}$ respectant les $E_{s}$ envoie $K^{(i)}$ dans
$A^{(i)}$, d'où une flèche $\psi : GK \to G(\widetilde{A}) \simeq A$ en sens
inverse. Le feuillet s'arrête après trois lignes, les deux tiers inférieurs
étant restés blancs.

\subsection*{Le facteur $(-1)^{i-1}(i-1)!$, et la réponse à la question}

La page 20 numérote trois propositions. Des deux premières il ne reste que le
mot et le numéro, la prose qui les énonçait étant perdue ; la troisième se
laisse reconstituer, non par son énoncé mais par sa démonstration, qui occupe
le haut de la page 21.

\textbf{Proposition (page 21).} \emph{Sur la composante de degré $i$,}
$\varphi \psi = (-1)^{i-1}(i-1)! \cdot \mathrm{id}$.

La démonstration est celle du feuillet, et elle est bonne : $CB \to CK$ et
$GB \to GK$ étant surjectifs, on se ramène à l'anneau universel $B$ ; pour
$P \in B^{(i)}$ on a $\varphi\psi$ appliqué à la classe de $P$ égal à
$\gamma^{i}(P - \varepsilon(P)) \bmod B^{(i+1)}$, et
\[
  \gamma^{i}\bigl(P - \varepsilon(P)\bigr) - (-1)^{i-1}(i-1)!\,P
  \;\in\; B^{(i+1)} .
\]
\footnote{La page 20 écrit ce coefficient $(-1)^{i}(i-1)!$ ; c'est
$(-1)^{i-1}(i-1)!$ qui est correct, comme la page 21 le démontre et comme le
vérifie le cas géométrique de la page 9, où $c^{r}(\xi^{r}) =
(-1)^{r-1}(r-1)!\,\eta^{r}$ tandis que $\operatorname{ch}_{r}(\xi^{r}) =
\eta^{r}$. La page 20 est un feuillet très serré dont la prose de liaison n'est
pas lisible ; la leçon de la page 21 est retenue.} Il en résulte que
$\varphi_{K}$ est toujours surjectif, et qu'il devient bijectif après
tensorisation par $\mathbb{Q}$ ; le feuillet énonce la surjectivité en
corollaire.

Il pose alors la question, en toutes lettres : « Est-ce que $\varphi_{K}$ est
toujours bijectif ??? », et répond « Non », le mot encadré, en renvoyant à
l'exemple qui suit.

\textbf{Exemple (page 21).} Soit $K = \mathbb{Z}[L]/(L-1)^{n}$, avec
$\lambda_{t}(L) = 1 + Lt$ --- c'est-à-dire $K = K(\mathbb{P}^{n-1})$, comme le
feuillet le note en marge. La filtration $\gamma$ y est la filtration par les
puissances de $(L-1)$, si bien que $GK = \mathbb{Z}[\zeta]/(\zeta^{n})$ est sans
torsion. Or $CK$ a de la torsion, engendrée par la relation
\[
  (1+\zeta)^{\,\cdot} \;=\; 1 + (-1)^{n-1}(n-1)!\,\zeta^{n} + \cdots
\]
Un homomorphisme surjectif d'un anneau avec torsion sur un anneau sans torsion
n'est pas injectif : $\varphi_{K}$ ne l'est donc pas.\footnote{Le générateur de
$CK$ est écrit $\xi$ sur le feuillet ; il est écrit $\zeta$ ici, la lettre
$\xi$ servant déjà à $1 - L^{-1}$ aux pages 9 à 13 et à une classe de $K(X)$
aux pages 42 et 43. L'exposant de $(1+\zeta)$ et le générateur exact de l'idéal
sont illisibles, ainsi qu'une note marginale en diagonale dont seuls l'idéal et
l'anneau $\mathbb{Z}[\zeta]/\zeta^{n}$ se laissent lire ; l'argument, lui, ne
dépend d'aucun des deux, et n'est pas reconstruit au-delà.} C'est le
contre-exemple des pages 9 à 13 --- le même facteur, la même torsion --- énoncé
sans géométrie.

\subsection*{L'anneau libre, et deux énoncés de finitude}

\textbf{Lemme (page 22).} Sur l'anneau libre $B = \mathbb{Z}\bigl[
(\lambda^{i}(X_{s}))_{1 \leqslant i \leqslant \varepsilon(s)},\,
(\lambda^{d}(Y_{t}))_{d \geqslant 1}\bigr]$, engendré par une famille
géométrique $(X_{s})$ et une famille libre $(Y_{t})$, et muni de la filtration
$B^{(k)}$ par les monômes de poids $\geqslant k$ en les $\gamma^{i}(X_{s} -
\varepsilon(X_{s}))$ et $\gamma^{d}(Y_{t} - \varepsilon(Y_{t}))$, la classe
totale $x \mapsto \widetilde{c}(x)$ vérifie
\[
  \widetilde{c}(x+y) = \widetilde{c}(x)\,\widetilde{c}(y),
  \qquad
  c^{i}(xy) = Q_{i}\bigl(c^{1}(x), \ldots ;\, c^{1}(y), \ldots\bigr).
\]
La première résulte de $\gamma_{t}(x' + y') = \gamma_{t}(x')\gamma_{t}(y')$
pour $x' = x - \varepsilon(x)$ ; la seconde de la formule correspondante pour
les $\gamma^{i}$.\footnote{Le feuillet écrit la seconde ligne avec
$\widetilde{c}(xy)$ au membre de gauche, après avoir biffé une ligne qui
portait $c^{i}(xy)$. C'est la ligne biffée qui est correcte : le membre de
droite est homogène de degré $i$. La leçon rayée est retenue ici.} Le feuillet
annonce ensuite que $\gamma^{i}(xy)$ est un polynôme universel en les
$\gamma^{k}(x)$ et $\gamma^{h}(y)$ de poids $\geqslant i$, et s'arrête sans
l'écrire ; c'est la première des choses que le dossier promet et ne donne pas.

\textbf{Proposition (page 24).} \emph{Soit $x$ avec $\varepsilon(x) = n$ et
$\lambda_{t}(x)$ polynomial de degré $\leqslant n$. Alors
$\gamma^{i}(x - n) = 0$ pour $i > n$.}

C'est immédiat : $\gamma_{t}(x - n) = \lambda_{t/(1-t)}(x)\,(1-t)^{n} =
\sum_{i \leqslant n} \lambda^{i}(x)\, t^{i}(1-t)^{n-i}$, polynôme de degré
$\leqslant n$ en $t$.\footnote{Le feuillet donne un argument différent, passant
par la nilpotence, dont la moitié des mots manque ; la démonstration ci-dessus
est celle que la formule elle-même fournit. La \emph{Remarque} et le
\emph{Corollaire} qui suivent sur le même feuillet portent sur la condition
exacte de nullité de $n!\,x^{n+1}$ et sur la filtration de $\gamma^{m}(X - m)$
dans l'anneau $\mathbb{Z}[(\lambda^{i}(X))_{i}]$ ; ils sont trop lacunaires pour
être restitués, et rien n'en est reconstruit ici.} L'énoncé dit qu'un élément de
rang fini a une filtration $\gamma$ bornée par son rang : c'est la propriété
qui rend la filtration utilisable en géométrie, où les fibrés ont un rang.

\pagerange{32}{37}
\section*{Construction des Chern covariants, et $K^{\bullet}$ et $A^{\bullet}$ à supports (pages 32 à 37)}

Cette course est la plus tardive du dossier et la plus lacunaire : elle est
écrite au crayon puis à l'encre sur les versos d'un texte dactylographié, et la
prose qui relie ses formules est presque entièrement perdue aux pages 36 et 37.
Ce qui suit s'en tient à ce que portent les formules et les titres.

\subsection*{Les catégories dérivées de complexes parfaits}

Soit $f : X \to Y$ localement de type fini. On note $D_{\mathrm{parf}}(f)$ la
catégorie dérivée des complexes de $\mathcal{O}_{X}$-Modules parfaits
relativement à $f$, et $K_{\mathrm{parf}}(f)$ le groupe de Grothendieck
correspondant.\footnote{Le feuillet qualifie cette catégorie d'un adjectif lu
« hégélienne », que la transcription donne comme incertain et que rien n'éclaire.
Il n'est pas repris ici. Le mot « parfait » figure deux fois sur la ligne, la
première biffée.} Pour $Z \subset X$ fermé, on définit de même
$D_{\mathrm{parf},Z}(X)$ et $K_{Z}(X)$, et pour $\Phi$ une famille de supports
\[
  D_{\mathrm{parf},\Phi}(X) = \varinjlim_{Z \in \Phi} D_{\mathrm{parf},Z}(X),
  \qquad
  K_{\Phi}(X) = \varinjlim_{Z \in \Phi} K_{Z}(X) .
\]
Pour $i : Z \to X$ on a $\mathrm{R}\,i_{*}$ vers $D_{Z}(X)$ et $i_{!}$ vers
$K_{Z}(X)$ ; et si $i_{n} : Z_{n} \to X$ parcourt les voisinages
infinitésimaux de $Z$ dans $X$, le foncteur
\[
  \varinjlim_{n} D_{\mathrm{parf}}(i_{n}) \longrightarrow D_{Z}(X)
\]
est pleinement fidèle.\footnote{L'objet dont $\mathrm{R}i_{*}$ est pris et le
qualificatif des voisinages $Z_{n}$ sont illisibles. L'énoncé donné est celui
que la flèche écrite supporte, et pas davantage.}

\subsection*{Les groupes de Chow à supports, et les classes de Chern}

La page 34 demande, en parallèle de $K_{\Phi}(X)$, un anneau de Chow à supports
$A_{\Phi}(X)$ --- qui n'a pas d'unité, l'élément $1$ n'ayant pas de support
dans $\Phi$, d'où la parenthèse « (pseudo-)anneaux \ldots{} par voie d'éliminer
$1$ » --- et des classes de Chern
\[
  c^{i}_{\Phi} : K_{\Phi}(X) \longrightarrow A_{\Phi}(X), \qquad i > 0,
\]
soumises à la formule de Whitney écrite sans unité :
\[
  c^{n}_{\Phi}(x+y) \;=\; c^{n}_{\Phi}(x) + c^{n}_{\Phi}(y)
  \;+\; \sum_{\substack{p+q=n \\ p,\,q > 0}} c^{p}_{\Phi}(x)\,c^{q}_{\Phi}(y),
\]
et compatibles avec les homomorphismes d'oubli du support :

\begin{tikzcd}[column sep=large, row sep=small]
K_{\Phi}(X) \arrow[r] \arrow[d, "c_{\Phi}"'] & K(X) \arrow[d, "c"] \\
A_{\Phi}(X) \arrow[r] & A(X)
\end{tikzcd}

Le feuillet ne construit rien de tout cela : il pose trois cas, et les trois
sont des questions. Pour $X$ lisse, prendre les cycles à support dans $\Phi$.
Pour $X$ général, se demander si la « $\lambda$-procédure » --- c'est-à-dire la
définition de $A$ comme gradué d'une filtration $\gamma$ --- s'applique à
$K_{\Phi}(X)$. Et enfin : des classes de Chern cohomologiques ? Aucune des
trois n'est traitée.\footnote{Le feuillet écrit, au milieu de cette page, en
anglais : « l'anneau de Chow \emph{ordinaire} whatever that means ». C'est le
seul endroit du dossier où il change de langue.}

\subsection*{Riemann-Roch comme mesure d'un défaut}

La page 36 est un palimpseste dont presque toute la prose est perdue. Ce que
ses formules dessinent est le programme suivant. On veut des groupes
$K^{\Phi}(X/S)$ et $A^{\Phi}(X/S)$ \emph{covariants} en $X$ au-dessus de $S$,
et entre eux un caractère de Chern qui commute strictement aux images
directes :
\[
  \operatorname{ch}^{Y/S}\bigl(g_{!}(x)\bigr)
  \;=\; g_{*}\bigl(\operatorname{ch}^{X/S}(x)\bigr) .
\]
Le caractère de Chern ordinaire, lui, ne commute pas aux images directes ---
c'est précisément l'objet du théorème. La phrase qui subsiste, « Mais
« R.R.» s'interprète comme \ldots{} », dit donc
que Riemann-Roch est ce qui traduit l'écart entre les deux : la classe de Todd
est la comparaison entre le caractère de Chern covariant et le caractère de
Chern contravariant.\footnote{Cette lecture est de l'éditeur. Les mots qui
relient les formules de la page 36 sont perdus au point que la phrase citée
n'est lisible qu'à un tiers ; ce qui est avancé ici est ce que la disposition
des formules impose, et rien n'est reconstruit au-delà. C'est, en termes
d'aujourd'hui, la construction d'une transformation naturelle de la $K$-théorie
vers l'homologie commutant aux images directes propres --- ce que le théorème
de Riemann-Roch singulier fera vingt ans plus tard.}

La page 37 dresse une échelle à quatre montants comparant les deux variances :

\begin{tikzcd}[column sep=small, row sep=small, nodes={font=\scriptsize}]
K^{\bullet}(i_{n}) \arrow[r, "u"] \arrow[d, "\alpha"'] & \varinjlim_{n} K^{\bullet}(i_{n}) \arrow[r] \arrow[d, "\beta"] & K^{\bullet}_{X}(X') \arrow[r] \arrow[d, "\gamma"] & K^{\bullet}(X') \arrow[d, "\delta"] \\
K_{\bullet}(X_{n}) \arrow[r, "u'"'] & \varinjlim_{n} K_{\bullet}(X_{n}) \arrow[r] & K^{X}_{\bullet}(X') \arrow[r, "w'"'] & K_{\bullet}(X')
\end{tikzcd}

Les quatre flèches verticales portent la même glose, et c'est la seule chose
que la page dise d'elles.\footnote{La glose est lue « isom. si $X'/X$ rig. ».
Ce sont les comparaisons entre la $K$-théorie des fibrés et celle des faisceaux
cohérents, dont on sait qu'elles sont bijectives sous une hypothèse de
régularité ; l'abréviation n'est pas tranchée par la transcription et n'est pas
interprétée ici. Un bloc de trois lignes au milieu du feuillet est entièrement
raturé, dont ne subsistent en marge que $w'v' = \mathrm{id}$ et $v'w' =
\mathrm{id}$.} Suit l'expression $\operatorname{ch}(i_{!}x)\cdot
\operatorname{Todd}(X'/S)$, qui est le membre de gauche de Riemann-Roch relatif,
puis un titre --- « Filtration sur $K^{\bullet}_{X}(X')$ » --- sous lequel rien
n'est écrit.

\pagerange{40}{41}
\section*{R.R. des surfaces : le carré, les trois formules, et la fibration (pages 40 et 41)}

La dernière course du dossier, annoncée par la chemise « R.R. des surfaces, et
autres cas particuliers divers », est une suite d'exemples calculés. Elle
n'emploie que l'appareil classique : caractère de Chern, classe de Todd,
$f_{!}$ et $f_{*}$. Aucun complexe, aucun support, aucune catégorie dérivée.

\subsection*{Le carré}

Soit $f : X \to Y$ un morphisme propre de variétés non
singulières.\footnote{La page 40 ne dit que « morphisme propre ». La
non-singularité est indispensable pour que $\operatorname{Todd}(X)$ et
$\operatorname{Todd}(Y)$ aient un sens et pour que $f_{!}$ opère sur $K(X) =
K^{0}(X)$ ; elle est ajoutée ici, et les exemples des pages 42 à 44 la
supposent tous.}

\begin{tikzcd}[column sep=large, row sep=large]
K(X) \arrow[r, "f_{!}"] \arrow[d, "\operatorname{ch}\,\cdot\,\operatorname{Todd}(X)"'] & K(Y) \arrow[d, "\operatorname{ch}\,\cdot\,\operatorname{Todd}(Y)"] \\
A(X) \otimes \mathbb{Q} \arrow[r, "f_{*}"'] & A(Y) \otimes \mathbb{Q}
\end{tikzcd}

\begin{equation*}
  \operatorname{Todd}(Y) \cdot \operatorname{ch}\bigl(f_{!}(\xi)\bigr)
  \;=\; f_{*}\bigl(\operatorname{Todd}(X) \cdot \operatorname{ch}(\xi)\bigr)
  \tag{1}
\end{equation*}

C'est le théorème de Grothendieck-Riemann-Roch.\footnote{Le feuillet dessine le
carré en quatre lignes plutôt qu'en deux, décomposant chaque flèche verticale
en trois : $\widetilde{c}$ vers $\widetilde{A}$, puis $h$ vers $A \otimes
\mathbb{Q}$, puis $\tau_{h}$. Le composé est le produit du caractère de Chern
par la classe de Todd, et c'est lui qui est dessiné ici ; la décomposition du
feuillet dit d'où vient chacun des deux facteurs.}

\subsection*{Les trois variantes}

\emph{Pour $f$ une immersion fermée}, la formule de projection permet de
rentrer $\operatorname{Todd}(Y)^{-1}$ sous $f_{*}$ et les classes de Todd de
$X$ et de $Y$ se combinent en celle du fibré normal $N$ de $X$ dans $Y$ :
\begin{equation*}
  \operatorname{ch}\bigl(f_{!}(\xi)\bigr)
  \;=\; f_{*}\Bigl(\operatorname{ch}(\xi)\cdot \operatorname{Todd}(N)^{-1}\Bigr)
  \tag{1 bis}
\end{equation*}
\footnote{Le feuillet écrit $T_{Y/X}$ là où l'on écrit aujourd'hui $N_{X/Y}$ :
c'est bien le fibré normal, et non un fibré tangent relatif, comme le montre
l'exposant $-1$ et comme le confirme la formule (2) de la même page, où
$\lambda_{-1}$ lui est appliqué. La lettre du feuillet est conservée dans la
transcription et traduite ici.}

Il annonce ensuite pour ce cas « une formule meilleure », numérotée (2), qui
donnerait la classe de Chern de degré $q+d$ de $f_{!}(\xi)$ à partir de
$\widetilde{c}(\xi) \ast \lambda_{-1}(\widetilde{c}(N))$ et de $c^{q}(N)$ : ce
serait la version entière du théorème, ce que l'on appellera Riemann-Roch sans
dénominateurs, et c'est la suite directe du calcul des pages 9 à 13. Sa forme
exacte n'est pas établie ici.\footnote{La transcription signale que la ligne
inférieure de (2), à l'encre bleue et portée sous une barre, est lue comme le
dénominateur d'une fraction mais peut tout aussi bien être une variante
ajoutée après coup. Les deux lectures donnent des énoncés différents, et rien
sur le feuillet ne permet de choisir. L'énoncé n'est donc pas donné ici, ni
sous une forme ni sous l'autre. Suivent sur le feuillet une référence
« (5.13) » et la formule $f_{!}(\tau_{X}(F)) = \tau_{Y}(F)$, dont le contexte
est perdu.}

\emph{Pour $f$ une fibration}, avec $E$ un fibré vectoriel sur $X$ et
$E^{(i)} = \mathrm{R}^{i}f_{*}(E)$, on a $f_{!}(E) = \sum_{i}(-1)^{i}E^{(i)}$
et donc
\begin{equation*}
  \operatorname{Todd}(Y) \sum_{i}(-1)^{i}\operatorname{ch}\bigl(E^{(i)}\bigr)
  \;=\; f_{*}\bigl(\operatorname{Todd}(X)\cdot\operatorname{ch}(E)\bigr)
  \tag{1 ter}
\end{equation*}

\subsection*{Ce que (1 ter) contient de plus que Hirzebruch}

C'est le meilleur passage du dossier, et il tient en deux observations.

\emph{En degré $0$}, (1 ter) redonne Riemann-Roch sur la fibre.\footnote{Le
feuillet écrit « Prenant les \ldots{} des deux membres », le mot gouvernant
toute la réduction étant illisible. Ce que le calcul exige est : les
composantes de degré $0$. C'est ainsi qu'il est restitué ici, et la lacune est
signalée.} Le membre de gauche devient $\sum_{i}(-1)^{i}\operatorname{rg}
E^{(i)} = \chi(Z;\, E|_{Z})$ pour $Z$ une fibre. Le membre de droite devient le
degré de la composante de degré $p = \dim Z$ de $i^{*}\bigl[
\operatorname{Todd}(X)\operatorname{ch}(E)\bigr]$, et l'on calcule
\[
  i^{*}\bigl(\operatorname{Todd}(X)\bigr)\cdot i^{*}\bigl(\operatorname{ch}(E)\bigr)
  \;=\; \operatorname{Todd}\bigl(i^{*}T_{X}\bigr)\operatorname{ch}\bigl(E|_{Z}\bigr)
  \;=\; \operatorname{Todd}(T_{Z})\operatorname{ch}\bigl(E|_{Z}\bigr) ,
\]
la dernière égalité parce que le fibré normal d'une fibre est
trivial.\footnote{C'est la seule chose que la réduction demande, et elle vaut
dès que $f$ est lisse : $N_{Z/X} \simeq f^{*}T_{Y,y}$ est trivial, donc
$i^{*}T_{X}$ et $T_{Z}$ ont même classe de Todd. La lissité n'est pas énoncée
sur le feuillet et est fournie ici.} On retrouve la formule de Hirzebruch pour
la fibre.

\emph{En degré $\geqslant 1$}, en revanche, (1 ter) donne les classes de Chern
de $\sum_{i}(-1)^{i}E^{(i)}$ --- et cela, comme le feuillet le souligne, n'est
pas contenu dans la formule de Hirzebruch proprement dite. Il en tire un
énoncé qui n'est pas trivial : \emph{pour $n$ grand, la classe de Chern de
$f_{*}(E(n))$ dépend polynomialement de $n$}.\footnote{Le feuillet porte cet
énoncé au bas de la page 41 sous une longue barre terminée par une flèche, et
sa dernière ligne se poursuit en tête de la page 42 : c'est une note ajoutée au
résultat, non la suite du texte. Les mots « pour $n$ grand » et l'adjectif
« modifiée » qualifiant la classe de Chern sont donnés comme incertains par la
transcription ; « modifiée » n'est pas interprété ici. Une note marginale sur
le bord inférieur gauche, dont ne se laissent lire que « ensuite » et « dans le
cas général », n'est pas restituée.}

\pagerange{42}{43}
\section*{Le développement en $n$, et une courbe dans une courbe (pages 42 et 43)}

\subsection*{La dépendance polynomiale en $n$}

Si $h$ est la première classe de Chern du faisceau tordant et $E(n) = E \otimes
\mathcal{O}(n)$, alors $\operatorname{ch}(E(n)) = \operatorname{ch}(E)e^{nh}$,
d'où
\[
  f_{*}\bigl(\operatorname{Todd}(X)\operatorname{ch}(E(n))\bigr)
  \;=\; \sum_{k} n^{k}\, f_{*}\!\left(\frac{h^{k}}{k!}\,
  \operatorname{Todd}(X)\,\operatorname{ch}(E)\right),
\]
somme finie puisque $h$ est nilpotente. C'est la démonstration de l'énoncé
annoncé au bas de la page 41 : le membre de droite est un polynôme en $n$ à
coefficients indépendants de $n$.\footnote{Le feuillet écrit $\xi$ pour cette
classe, sur le feuillet même où $\xi$ désigne déjà l'élément $(n,d)$ de $K(X)$.
Elle est écrite $h$ ici.}

\subsection*{Morphismes propres particuliers}

Le feuillet ouvre alors, sous un titre de sa main séparé par un trait ondulé,
une suite de cas particuliers. Le premier est celui d'un morphisme d'une courbe
dans une variété.

Dans tout ce qui suit, $X$ est la courbe source et $Y$ le but.\footnote{C'est ce
qu'exigent toutes les formules de ces trois feuillets, et ce que porte la
formule encadrée de la page 42, « $f : X \to Y$ ». Les lettres vacillent
pourtant : le titre du feuillet 42 annonce « morphismes d'une courbe dans une
variété $X$ », et la lettre est effacée sous un pâté à chacune de ses premières
occurrences, ce qui donne à penser qu'il a renommé source et but en cours de
route. Rien n'a été réparé dans la transcription ; ici, les lettres sont
tenues, et chaque endroit où le feuillet porte l'autre est signalé.}

Sur une courbe non singulière, $K(X) \simeq \mathbb{Z} + \operatorname{Pic}(X)$,
l'isomorphisme étant donné par la classe de Chern. Pour $\xi = (n,d)$ :
\[
  \operatorname{ch}(\xi) = n + d,
  \qquad c(T_{X}) = 1 - k_{X},
  \qquad \operatorname{Todd}(X) = \frac{-k_{X}}{1 - e^{k_{X}}}
  = 1 - \tfrac{1}{2}k_{X},
\]
$k_{X}$ étant la classe canonique, et donc
\[
  \operatorname{ch}(\xi)\operatorname{Todd}(X)
  \;=\; \bigl(1 - \tfrac{1}{2}k_{X}\bigr)(n+d)
  \;=\; n + \bigl(d - \tfrac{1}{2}n\,k_{X}\bigr).
\]
On reconnaît la formule (1) de la page 3, écrite trente-neuf feuillets plus
tôt : la classe de Todd d'une courbe est $1 - \tfrac12 k_{X}$ dans les deux
cas, et le degré de ce produit est la caractéristique d'Euler-Poincaré.

\subsection*{a) Le but est une courbe}

Soit $k_{Y}$ son diviseur canonique, $\operatorname{Todd}(Y) = 1 -
\tfrac12 k_{Y}$, et $\delta$ le degré de $f$, de sorte que $f_{*}(1) = \delta$.
Alors
\[
  f_{*}\bigl(\operatorname{ch}(\xi)\operatorname{Todd}(X)\bigr)
  \;=\; n\delta + \bigl(f_{*}(d) - \tfrac{1}{2}n\,f_{*}(k_{X})\bigr),
\]
et en divisant par $\operatorname{Todd}(Y)$ :
\[
  \operatorname{ch}\bigl(f_{!}(\xi)\bigr)
  \;=\; \bigl(1 + \tfrac{1}{2}k_{Y}\bigr)
  \Bigl( n\delta + f_{*}(d) - \tfrac{1}{2}n\,f_{*}(k_{X}) \Bigr)
  \;=\; n\delta + \Bigl[\tfrac{1}{2}n\delta\,k_{Y} + f_{*}(d)
  - \tfrac{1}{2}n\,f_{*}(k_{X})\Bigr],
\]
soit, dans les coordonnées $(\text{rang},\, c^{1})$ :
\[
  f_{!}(\xi) \;=\; \Bigl(n\delta,\ \ \tfrac{1}{2}n\delta\,k_{Y}
  + f_{*}(d) - \tfrac{1}{2}n\,f_{*}(k_{X})\Bigr),
  \qquad f : X \to Y .
\]
\footnote{Le feuillet écrit le membre de gauche de l'avant-dernière ligne
$\varphi_{X}(f_{!}(\xi))$ : c'est $\varphi_{Y}$ qu'il faut, puisque $f_{!}(\xi)$
vit dans $K(Y)$. La ligne suivante de la page 43 écrit d'ailleurs $\varphi_{Y}$.
La lettre du feuillet est reproduite telle quelle dans la transcription, et
corrigée ici. Dans la formule encadrée il distingue les deux classes canoniques
en portant la base en indice, $k_{y}$ et $k_{x}$, là où les lignes précédentes
portaient $k'$ et $k$ ; l'indexation par la base est ce qui est tenu ici.}

\subsection*{Le revêtement}

Appliquant cela à $\xi = [\mathcal{O}_{X}]$, c'est-à-dire $n = 1$, $d = 0$,
pour $f$ un revêtement de degré $\delta$ : le faisceau $f_{*}(\mathcal{O}_{X})$
est un fibré vectoriel \emph{sur $Y$}, de rang $\delta$, et sa première classe
de Chern vaut
\[
  c^{1}\bigl(f_{*}\mathcal{O}_{X}\bigr)
  \;=\; \tfrac{1}{2}\bigl(\delta\,k_{Y} - f_{*}(k_{X})\bigr).
\]
\footnote{Le feuillet place $f_{*}(\mathcal{O}_{X})$ « sur $X$ » là où le calcul
le demande sur $Y$ ; la lettre est reproduite telle quelle dans la
transcription et corrigée ici. Il indique aussi un degré, dont la transcription
ne retient que le fragment « $-\,g\,\delta$ » ; le degré que la formule donne
est $\deg c^{1}(f_{*}\mathcal{O}_{X}) = \delta(g_{Y}-1) - (g_{X}-1)$, ce qui
est bien $\chi(\mathcal{O}_{X}) - \delta\,\chi(\mathcal{O}_{Y})$ comme il se
doit. Le fragment lisible ne suffit pas à dire si c'est ce qu'il avait écrit, et
rien de plus n'est affirmé ici.}

\pagerange{43}{44}
\section*{b) Une courbe dans une surface, et la formule d'adjonction (pages 43 et 44)}

Soit maintenant $Y$ une surface non singulière et $f : X \to Y$ un morphisme
depuis une courbe, d'image la courbe $C$, de degré $\delta$ sur son
image.\footnote{Les trois lignes qui ouvrent ce cas sont barrées d'un long
trait et portent deux additions interlinéaires ; seuls « classe canonique » et
deux $k_{Y}$ s'en laissent lire. Le feuillet note $k_{Y}^{1}$ et $k_{Y}^{2}$
les deux classes de Chern du fibré tangent, la première étant la classe
canonique ; elles sont écrites $k_{Y}$ et $c^{2}(T_{Y})$ ici.}

De $c(T_{Y}) = 1 - k_{Y} + c^{2}(T_{Y})$ on tire
\[
  \operatorname{Todd}(Y) = 1 - \tfrac{1}{2}k_{Y}
  + \tfrac{1}{12}\bigl[k_{Y}^{2} + c^{2}(T_{Y})\bigr],
  \qquad
  \operatorname{Todd}(Y)^{-1} = 1 + \tfrac{1}{2}k_{Y}
  + \tfrac{1}{12}\bigl[k_{Y}^{2} - c^{2}(T_{Y})\bigr] \cdot
\]
\footnote{Le signe qui joint les deux termes du crochet de
$\operatorname{Todd}(Y)^{-1}$ est recouvert sur le feuillet et la transcription
ne le tranche pas. Il ne sert nulle part dans la suite : le calcul qui vient
n'utilise que les termes de degré $\leqslant 1$, puisque $f_{*}$ d'une classe
de degré $\leqslant 1$ sur une courbe ne dépasse pas le degré $2$ sur la
surface. Le terme de degré $2$ est écrit ici sous la forme qu'il a, sans que
rien n'en dépende.}

On a $f_{*}(1) = \delta\,C$, et $f_{*}(k_{X})$, $f_{*}(d)$ sont des classes de
$0$-cycles. D'où, pour $\xi = (n,d)$ :
\[
  \operatorname{ch}\bigl(f_{!}(\xi)\bigr)
  \;=\; n\delta\,C + \Bigl[\,f_{*}(d) - \tfrac{1}{2}n\,f_{*}(k_{X})
  + \tfrac{1}{2}n\delta\,(C\cdot k_{Y})\,\Bigr],
\]
et en passant à la classe de Chern totale,
\[
  c\bigl(f_{!}(\xi)\bigr) \;=\; 1 + n\delta\,C
  + \Bigl[\,\tfrac{1}{2}n^{2}\delta^{2}C^{2} - \tfrac{1}{2}n\delta\,(C\cdot k_{Y})
  - f_{*}(d) + \tfrac{1}{2}n\,f_{*}(k_{X})\,\Bigr].
\]
Le passage se fait par $c = \exp\bigl(\operatorname{ch}_{1} -
\operatorname{ch}_{2}\bigr)$ à l'ordre $2$, ce que le feuillet écrit
exactement.\footnote{Et c'est le quatrième visage du nombre du dossier : les
deux premiers termes de $\log c = \sum_{m}(-1)^{m-1}(m-1)!\operatorname{ch}_{m}$
sont $\operatorname{ch}_{1} - \operatorname{ch}_{2}$. Le feuillet écrit
l'exponentielle et ne nomme pas l'identité ; le rapprochement avec le $(i-1)!$
de la page 21 et le $(r-1)!$ de la page 9 est de cette lecture.} Pour
$n = \delta = 1$ :
\[
  c\bigl(f_{!}(\xi)\bigr) \;=\; 1 + C
  + \Bigl[\,\tfrac{1}{2}\bigl(C^{2} - C\cdot k_{Y}\bigr) - f_{*}(d)
  + \tfrac{1}{2}f_{*}(k_{X})\,\Bigr].
\]

\subsection*{La dernière page, et la formule qu'elle n'écrit pas}

Le dernier feuillet applique ceci à $f : \widetilde{C} \to Y$, la normalisée
d'une courbe $C$ tracée sur la surface, avec $\xi = [\mathcal{O}_{\widetilde{C}}]$.
Il trouve
\[
  c^{2}\bigl(\mathcal{O}_{C}\bigr) \;=\; \tfrac{1}{2}
  \Bigl[\,C^{2} - C\cdot k_{Y} + f_{*}\bigl(k_{\widetilde{C}}\bigr)\,\Bigr],
\]
puis, $Y$ étant compacte et $\widetilde{C}$ de genre $g$, en prenant le degré
des deux membres :
\[
  \deg c^{2}\bigl(\mathcal{O}_{C}\bigr)
  \;=\; \tfrac{1}{2}\bigl[\,C^{2} - C\cdot k_{Y} + 2g - 2\,\bigr]
  \;=\; \tfrac{1}{2}\bigl(C^{2} - C\cdot k_{Y}\bigr) + g - 1 .
\]
\footnote{La surface est appelée $Y$ dans les formules et $X$ dans la phrase qui
les commente, sans que rien sur le feuillet distingue les deux ; l'indice du
dernier $k$ est recouvert d'encre et se lit sur la ligne suivante, et le but de
$f : \widetilde{C} \to \,\cdot$ est effacé sous un pâté. Ce que la cohérence
demande est ce qui est écrit ici : $Y$ est la surface, $f$ va de $\widetilde{C}$
dans $Y$, et l'opérateur $K_{2}$ du feuillet est le $\deg$ de la page 3 et le
$K_{p}$ de la page 41.}

Le feuillet s'arrête là, et le reste de la page est blanc : c'est la fin du
dossier. Or il est à une identification de la formule d'adjonction. La classe
$[\mathcal{O}_{C}]$ vaut $1 - [\mathcal{O}_{Y}(-C)]$ dans $K(Y)$, d'où
$c(\mathcal{O}_{C}) = (1-C)^{-1}$ et $c^{2}(\mathcal{O}_{C}) = C^{2}$ lorsque
$C$ est lisse. Égaler les deux valeurs donne
\[
  C^{2} \;=\; \tfrac{1}{2}\bigl(C^{2} - C\cdot k_{Y}\bigr) + g - 1,
  \qquad \text{c'est-à-dire} \qquad
  2g - 2 \;=\; C^{2} + C\cdot k_{Y} .
\]
C'est la formule d'adjonction, obtenue comme corollaire de
Grothendieck-Riemann-Roch pour la normalisée d'une courbe sur une
surface.\footnote{Ce dernier pas n'est pas sur le feuillet. Il calcule le
premier membre et s'arrête ; le second calcul, immédiat, n'est nulle part dans
le dossier, et le rapprochement est de cette lecture. On notera qu'écrite dans
ce sens, la formule vaut pour $C$ lisse ; pour $C$ singulière, c'est la version
de gauche qui reste vraie et qui définit le genre géométrique de la
normalisée --- ce qui est sans doute pourquoi il l'a écrite dans ce sens-là.}

\end{document}
